\chapter{线程、调度与亲和性}

线程是操作系统调度的执行实体。多线程程序能利用多核心，但也引入调度、同步、内存可见性和资源竞争。理解线程，必须同时看用户态代码和内核调度。

\section{从任务到线程}

创建八个线程不等于拥有八个物理核心。硬件可能有 SMT，一个物理核心暴露两个逻辑 CPU；系统里还运行着其他进程；线程可能被迁移到不同核心。线程迁移会破坏缓存热度，过多线程会增加上下文切换。

CPU-bound 程序通常不应远多于核心数，I/O-bound 程序可以有更多并发等待。判断依据不是经验数字，而是运行时是否在计算、等待锁、等待 I/O 或被调度器抢占。

\topic{用户线程和内核线程}

C++ 的 \code{std::thread} 通常映射到操作系统线程。操作系统调度的是内核可见的执行实体，它拥有寄存器上下文、栈、调度状态和优先级。用户态协程则不同，协程切换通常不进入内核，它把等待点显式暴露给运行时。协程适合组织大量 I/O 等待，不会自动让 CPU 密集代码并行。

这一区分很重要。把 CPU 密集任务放进单线程协程事件循环，不会使用多个核心；把大量 I/O 等待都变成内核线程，可能造成过多栈内存、调度成本和上下文切换。运行时设计要把计算并行和等待并发分开。

\section{调度、阻塞和上下文切换}

上下文切换需要保存和恢复执行状态，可能破坏 cache 和 TLB 局部性。频繁阻塞、线程过多、锁竞争和系统调用都可能增加切换。性能计数器里的 context-switches 和 migrations 可以帮助定位这类问题。

Linux CFS 调度器大体追求公平，它会根据虚拟运行时间选择任务。实时调度策略和普通调度策略不同，容器和 cgroup 也会改变 CPU 配额。benchmark 时若机器上有其他负载，或者进程受 cgroup 限制，结果会明显波动。严谨实验应记录负载、CPU governor、容器限制和是否独占机器。

调度器不是只在“线程很多”时才重要。即使线程数等于核心数，线程也可能因为缺页、锁等待、I/O、信号、中断、抢占和 cgroup 配额离开 CPU。对于低延迟系统，一次迁核可能让热缓存失效；对于吞吐系统，频繁上下文切换会减少有效计算时间。性能报告如果只写线程数，不写线程是否真正持续运行在 CPU 上，就缺少关键证据。

Linux 调度源码阅读不建议一开始追完整策略。更好的入口是从一个用户态现象出发：为什么 \cmd{perf stat} 里 context-switches 增加，为什么线程会 migration，为什么绑定 CPU 后波动变小，为什么 cgroup 限额下吞吐呈周期性抖动。围绕这些问题，再去看调度实体、runqueue、唤醒路径和上下文切换路径，读出的内容才有边界。

\topic{阻塞、睡眠和自旋}

线程等待锁时，可以自旋，也可以睡眠。自旋避免进入内核，适合等待时间极短的场景；睡眠释放 CPU，适合等待时间较长的场景。很多锁实现会先短暂自旋，再进入 futex 等内核等待路径。这样低竞争时快，高竞争时不至于烧满 CPU。

选择自旋还是阻塞，本质是在 CPU 时间和唤醒延迟之间取舍。如果临界区很短且核心充足，自旋可能好；如果锁持有者被抢占或等待 I/O，自旋会浪费大量周期。不要在没有测量的情况下把自旋当成高性能。

futex 是理解 Linux 用户态锁的重要入口。常见 mutex 在无竞争时尽量只走用户态原子操作；竞争严重时，线程通过 futex 进入内核睡眠，锁释放时再唤醒等待者。这样做把低竞争 fast path 和高竞争 blocking path 分开。读锁性能时，应区分用户态自旋成本、futex 睡眠成本、唤醒成本和被唤醒后重新竞争锁的成本。

\section{亲和性、拓扑和移动设备}

线程亲和性让线程倾向或固定在某些 CPU 上运行。它可以改善缓存局部性和 NUMA 放置，也可能降低调度器自由度。使用亲和性前要理解拓扑：物理核心、SMT 线程、NUMA 节点和共享缓存层级。

亲和性实验要有对照。只绑定线程不绑定内存，可能仍然访问远端页面；只绑定内存不绑定线程，线程迁移后也可能变成远端访问。合理实验至少比较未绑定、只绑线程、只绑内存、线程和内存都按节点绑定几种情况。若平台没有 NUMA，也应记录物理核心和 SMT 拓扑，避免把两个重负载线程绑到同一个物理核心的两个逻辑 CPU 上后误判算法慢。

\topic{线程生命周期}

频繁创建销毁线程成本高，也让任务调度不可控。线程池把线程生命周期和任务生命周期分离，适合大量短任务。但线程池不是万能的，如果任务太小，调度成本会超过计算成本；如果任务阻塞，线程池可能被耗尽。

\topic{线程局部存储}

线程局部存储让每个线程拥有独立对象，适合缓存临时缓冲、统计计数和避免共享写。它可以减少锁竞争，但也可能增加内存占用，并让跨线程汇总更复杂。线程池里使用线程局部对象时，还要注意任务迁移：同一个逻辑请求的不同阶段可能被不同 worker 执行，不能把请求状态误放到线程局部。

\section{线程池语义和配置}

线程池把固定 worker 暴露成任务执行资源，但它不会自动解决调度语义。CPU 密集任务应该限制并发度，避免超过核心数太多；阻塞 I/O 任务如果占住 worker，可能让计算任务饥饿；任务内部再提交子任务并等待，可能造成线程池死锁。设计线程池时，要明确任务是否允许阻塞、是否允许嵌套提交、是否支持 work helping、是否需要独立 I/O 池。

线程池还会改变线程局部数据的含义。线程局部缓冲属于 worker，不属于请求。请求如果跨多个任务执行，就不能假设后续阶段还能访问同一个线程局部对象。运行时应把“worker 本地缓存”和“请求上下文”分成两个概念。

\topic{从任务到线程的映射}

并发程序里最容易混淆的是任务和线程。任务是逻辑工作单元，例如处理一段数组、解析一个文件块、执行一个 reduce 分区、发送一个网络请求。线程是操作系统调度实体，它拥有栈和寄存器上下文，会被放入 runqueue，由调度器分配 CPU 时间。一个任务可以在线程上执行，一个线程可以执行很多任务；二者不是同一概念。

Compute Systems Engine 最初的 \code{parallel_sum} 每次调用都创建一组线程，任务和线程几乎一一对应。这种写法适合教学，因为边界清楚；但它不适合大量短任务，因为线程创建、栈分配、调度注册和 join 都有成本。线程池把线程生命周期延长，让任务通过队列进入 worker。这样任务变轻了，但又引入队列、等待、关闭和异常传播问题。每前进一步，性能和语义都会同时变化。

设计运行时时，要先回答三个问题。第一，任务是否 CPU 密集。如果是，worker 数通常不应远超物理核心，除非需要隐藏内存或 I/O 等待。第二，任务是否可能阻塞。如果任务会等待磁盘、网络、锁或 future，把它放进 CPU 池可能耗尽 worker。第三，任务之间是否有依赖。如果一个任务提交子任务后等待，而线程池没有 work helping 或足够 worker，可能死锁。线程池不是“并发万能容器”，它只是执行资源的组织方式。

\topic{线程创建成本和任务粒度}

任务粒度是并行运行时最基本的参数。粒度太大，负载不均衡，一个慢任务拖住整体；粒度太小，队列操作、调度、同步和计数器开销超过有效计算。归约任务如果每个 worker 处理上千万元素，线程创建成本相对很小；如果每个任务只处理几十个元素，线程池也可能被调度开销吞掉。

选择粒度不能只靠固定常数。不同机器核心数不同，缓存大小不同，内存带宽不同，任务内核成本也不同。一个合理策略是先用 coarse chunk 保证每个任务有足够计算量，再根据负载不均衡逐步拆小。对均匀数组求和，按连续大块切分足够；对图遍历或 key 倾斜聚合，任务成本不均匀，可能需要动态调度、工作窃取或分片负载估计。

实验上要画任务粒度曲线。固定总输入和线程数，改变 chunk size，记录吞吐、任务数、队列等待、worker 空闲时间和尾部完成时间。曲线通常会出现一个中间区域：粒度太小，调度成本高；粒度太大，负载不均衡；中间区域才是合理范围。这个实验比背“任务数等于线程数几倍”更可靠。

\topic{阻塞任务和线程池耗尽}

考虑一个线程池有四个 worker。每个任务启动后提交一个子任务，然后等待子任务完成。如果四个 worker 都在执行父任务并等待子任务，而子任务还在队列里没有 worker 执行，线程池就死锁了。这个问题不是 mutex 死锁，却同样来自等待关系形成环。

解决方向有几种。第一，禁止 worker 内同步等待同一线程池的子任务，把依赖交给任务图调度。第二，支持 work helping：worker 等待时不睡眠，而是继续执行队列里的其他任务。第三，分离 CPU 池和阻塞池，让阻塞 I/O 不占满计算 worker。第四，增加结构化并发语义，让任务的生命周期和父子关系明确。选择哪种方式取决于运行时目标，但必须在接口文档中说明，不能让调用者猜。

本章先建立问题，后续任务运行时章节会实现更完整的任务图。这里要记住：线程池接口不只是 \code{submit}，还包括等待、取消、关闭、异常和嵌套提交语义。一个只会把函数塞进队列的线程池，很容易在真实系统中出现停不下来、等不到、吞异常或死锁的问题。

\topic{亲和性和拓扑记录}

亲和性优化的前提是拓扑记录。报告里至少要写清逻辑 CPU 数、物理核心数、SMT 关系、NUMA 节点、共享 L1/L2/L3 的范围和内存节点。Linux 上可以用 \cmd{lscpu}、\cmd{numactl --hardware}、\filepath{/sys/devices/system/cpu/} 下的拓扑文件和工具输出记录。没有这些信息，绑定 CPU 的结论很难解释。

绑定策略要和工作负载匹配。CPU 密集向量计算通常先使用每个物理核心一个线程，再测试 SMT；内存延迟敏感但单线程经常等待的任务可能从 SMT 获益；共享数据强的任务可能受益于放在共享 LLC 内；NUMA 分片任务应让线程和内存位于同一节点；跨分片通信强的任务要考虑互连成本。亲和性不是为了把线程“钉死”就结束，而是为了把执行位置、数据位置和共享层级对齐。

亲和性也可能伤害性能。过度绑定会减少调度器处理系统噪声的自由度；把两个重负载线程绑到同一物理核心的 SMT sibling 会互抢执行资源；把线程绑到一个 NUMA 节点但数据在另一个节点，会稳定地产生远端访问；在手机或容器环境里，系统可能限制可用 CPU，强行绑定到不可用 CPU 会失败。实验报告必须记录绑定是否成功。

\topic{上下文切换的真实代价}

上下文切换不是一个固定纳秒数。直接成本包括保存和恢复寄存器状态、切换内核调度结构和更新运行队列；间接成本包括 cache 热度下降、TLB 上下文影响、分支预测状态变化和被迁移到不同核心后的数据重新加载。若切换发生在锁持有者身上，其他线程可能自旋或睡眠等待；若切换发生在低延迟请求路径上，尾延迟可能被放大。

测上下文切换要区分自愿和非自愿。自愿切换通常来自线程主动阻塞，例如等待条件变量、I/O 或 futex；非自愿切换通常来自时间片耗尽、抢占或更高优先级任务。两者解决方向不同。自愿切换多，可能要减少阻塞、批处理或改变等待策略；非自愿切换多，可能线程数过多、CPU 竞争严重或系统负载干扰。

在 \cmd{perf stat} 或系统工具中看到 context-switches 增加时，不要立刻下结论。一个 I/O 服务有上下文切换可能是正常的；一个 CPU-bound 热循环频繁切换就值得怀疑。要结合 CPU 利用率、线程数、锁等待、runqueue 长度和任务类型分析。

\topic{Linux 调度源码阅读路线}

Linux 调度器很复杂，不适合一开始从所有策略全局阅读。本册建议按现象驱动。若你想理解普通线程为什么被调度，可以从核心调度入口、调度类、运行队列和上下文切换路径读起；若想理解唤醒延迟，可以沿 futex 或条件变量唤醒路径看到任务如何重新进入 runqueue；若想理解迁移，可以阅读负载均衡和 CPU affinity 相关路径；若想理解 cgroup 限额，可以阅读 CPU controller 对可运行时间的限制。

报告要写具体问题。例如“在 16 个线程测试中，migrations 明显增加，绑定 CPU 后波动下降”，然后阅读调度迁移和亲和性相关代码。不要写“Linux 使用 CFS 调度器，所以性能如何”这种大而空的句子。源码阅读的价值在于把用户态现象和内核数据结构连接起来：task、runqueue、调度类、唤醒、抢占、迁移和上下文切换。

\topic{线程局部缓存和内存生命周期}

线程局部缓存可以减少分配和共享写，但它改变了对象生命周期。线程池中的 worker 会长期存在，线程局部对象也会长期存在。如果线程局部缓冲按最大请求扩容，处理一次大请求后可能长期占用大量内存。若任务之间复用线程局部状态，还要保证不会泄漏上一个请求的数据。高性能系统里，缓存和生命周期总是绑在一起。

线程局部统计同样有读取语义问题。每个 worker 写自己的计数很快，但读取总数需要遍历所有 worker 状态。若 worker 正在写，读者是否需要精确值？如果只用于指标展示，可以接受近似或延迟汇总；如果用于协议决策，例如限流额度，就需要更强同步。把所有统计都改成线程局部之前，必须说明读取方的语义需求。

\topic{案例：从每次创建线程到固定线程池}

下面用归约说明线程生命周期的演化。baseline 是每次调用 \code{parallel_sum} 都创建 \code{worker_count} 个线程。它适合少量大任务，但不适合频繁短任务。改进方向是固定线程池：程序启动时创建 worker，归约请求被拆成多个 chunk，提交到队列，worker 执行后写 partial，主线程等待所有 chunk 完成。

这个改进引入了新问题。队列需要同步，提交任务需要分配或复用任务对象，等待所有 chunk 完成需要计数器或 latch，异常需要传回调用者，线程池关闭时不能丢任务。如果允许多个归约请求同时进入线程池，partial 的所有权和生命周期要隔离。一个高质量线程池章节不能只给一段 \code{while(true)} worker 循环，还要讲清这些语义。

本章只要求读者理解演化方向。后续运行时章节会把线程池拆成接口、队列、worker、任务状态、关闭协议和指标。这里先建立判断：当任务足够大且调用次数少，每次创建线程可以接受；当任务短且频繁，线程池值得引入；当任务之间有依赖，普通 FIFO 线程池不够，需要任务图或 work stealing。


\topic{高密度主线补充：从失败推导机制}

线程章必须从任务为什么不能直接等同于线程讲起。这一轮补充专门解决两个问题：第一，避免把本章写成概念枚举；第二，让每个概念都能从把日志作业拆成多个可执行任务的失败中自然推出。本章的核心失败不是“代码不会写”，而是线程数量、任务粒度、阻塞和调度位置没有边界，导致过度并行、上下文切换和不可复现实验。所以正文的顺序必须始终保持为：先给出具体任务，再写朴素方案，再观察失败信号，再引入机制，最后给出实验和边界。
这样的写法比直接给定义慢一些，但它更接近工程学习的真实路径。真实系统很少把问题包装成选择题；它只会表现为吞吐上不去、CPU 利用率怪、结果偶尔错、队列越积越多、重启后状态对不上、线上延迟突然变长。读者要学会从这些现象反推机制，而不是看到术语才想起术语。

\topic{任务和线程}

先把问题收窄到一个可推导的场景：一个输入分片为什么不应该必然对应一个系统线程。在把日志作业拆成多个可执行任务里，这个问题不会以术语形式出现，而是以结果不稳定、吞吐不增长、延迟抖动或恢复困难的形式出现。如果一上来只背任务和线程的定义，读者很容易把它当成孤立知识点；更好的顺序是先看朴素方案为什么合理，再看它在真实约束下怎样失败。

朴素方案通常是：每个分片创建一个 thread。这个方案的吸引力在于它贴近单线程直觉，代码短，状态少，测试也容易写。但是它隐含了几个没有说出口的假设：输入总是干净，资源近似无限，执行不会交错，失败不会发生，硬件成本和源码结构大体一致。一旦这些假设被打破，线程创建和调度开销随分片数膨胀，错误传播和关闭也难控制。这时继续在原方案上加 if 或加锁，往往只会把真正的问题埋得更深。

任务和线程就是从这个失败里长出来的概念。它的核心模型可以这样理解：任务是工作单元，线程是执行资源，运行时负责映射。这个模型必须同时放在三个层次看。源码层要说明谁读谁写、谁拥有对象、哪个状态允许变化；运行时层要说明线程、队列、任务和 I/O 怎样排队；硬件或系统层要说明缓存、页、调度、文件系统或网络怎样参与成本。三层缺一层，解释都会变形：只有源码层会看不见隐藏成本，只有硬件层会丢失业务语义，只有运行时层又容易把正确性和性能混在一起。

观察方法不能停在“跑一下变快了”。这一点在本章尤其重要：比较每任务一线程和固定线程池。实验要有 reference，要固定输入规模，要记录环境，要把冷启动、预热、核心循环、提交和清理分开。如果工具权限不足，也要写清楚不能观察什么，而不是把猜测写成结论。系统教材的严谨性来自证据链：现象、假设、对照、指标、结论边界。

任务和线程也有边界。任务太大负载不均，太小调度成本高。工程判断不是把一个技术推到所有地方，而是知道它在哪些约束下成立。读者在本章应形成一种习惯：每学到一个机制，都要问它解决了什么失败、引入了什么新成本、需要什么测试、在什么输入或部署环境下会失效。

\topic{过度订阅}

先把问题收窄到一个可推导的场景：线程数超过可用执行资源会发生什么。在把日志作业拆成多个可执行任务里，这个问题不会以术语形式出现，而是以结果不稳定、吞吐不增长、延迟抖动或恢复困难的形式出现。如果一上来只背过度订阅的定义，读者很容易把它当成孤立知识点；更好的顺序是先看朴素方案为什么合理，再看它在真实约束下怎样失败。

朴素方案通常是：按输入分片数量创建线程。这个方案的吸引力在于它贴近单线程直觉，代码短，状态少，测试也容易写。但是它隐含了几个没有说出口的假设：输入总是干净，资源近似无限，执行不会交错，失败不会发生，硬件成本和源码结构大体一致。一旦这些假设被打破，上下文切换、缓存失效和调度延迟增加。这时继续在原方案上加 if 或加锁，往往只会把真正的问题埋得更深。

过度订阅就是从这个失败里长出来的概念。它的核心模型可以这样理解：过度订阅让 runnable 线程争夺有限 CPU。这个模型必须同时放在三个层次看。源码层要说明谁读谁写、谁拥有对象、哪个状态允许变化；运行时层要说明线程、队列、任务和 I/O 怎样排队；硬件或系统层要说明缓存、页、调度、文件系统或网络怎样参与成本。三层缺一层，解释都会变形：只有源码层会看不见隐藏成本，只有硬件层会丢失业务语义，只有运行时层又容易把正确性和性能混在一起。

观察方法不能停在“跑一下变快了”。这一点在本章尤其重要：记录 runnable 数、context switch 和吞吐曲线。实验要有 reference，要固定输入规模，要记录环境，要把冷启动、预热、核心循环、提交和清理分开。如果工具权限不足，也要写清楚不能观察什么，而不是把猜测写成结论。系统教材的严谨性来自证据链：现象、假设、对照、指标、结论边界。

过度订阅也有边界。I/O 阻塞任务可能需要不同策略，不能只按核心数固定。工程判断不是把一个技术推到所有地方，而是知道它在哪些约束下成立。读者在本章应形成一种习惯：每学到一个机制，都要问它解决了什么失败、引入了什么新成本、需要什么测试、在什么输入或部署环境下会失效。

\topic{阻塞}

先把问题收窄到一个可推导的场景：一个 worker 阻塞在 I/O 或条件变量上是否还算占用计算资源。在把日志作业拆成多个可执行任务里，这个问题不会以术语形式出现，而是以结果不稳定、吞吐不增长、延迟抖动或恢复困难的形式出现。如果一上来只背阻塞的定义，读者很容易把它当成孤立知识点；更好的顺序是先看朴素方案为什么合理，再看它在真实约束下怎样失败。

朴素方案通常是：把阻塞任务和计算任务放同一个池。这个方案的吸引力在于它贴近单线程直觉，代码短，状态少，测试也容易写。但是它隐含了几个没有说出口的假设：输入总是干净，资源近似无限，执行不会交错，失败不会发生，硬件成本和源码结构大体一致。一旦这些假设被打破，计算线程被等待占住，吞吐下降且排队难解释。这时继续在原方案上加 if 或加锁，往往只会把真正的问题埋得更深。

阻塞就是从这个失败里长出来的概念。它的核心模型可以这样理解：阻塞会改变线程池容量模型，需要隔离或异步化。这个模型必须同时放在三个层次看。源码层要说明谁读谁写、谁拥有对象、哪个状态允许变化；运行时层要说明线程、队列、任务和 I/O 怎样排队；硬件或系统层要说明缓存、页、调度、文件系统或网络怎样参与成本。三层缺一层，解释都会变形：只有源码层会看不见隐藏成本，只有硬件层会丢失业务语义，只有运行时层又容易把正确性和性能混在一起。

观察方法不能停在“跑一下变快了”。这一点在本章尤其重要：记录线程状态、队列长度和等待原因。实验要有 reference，要固定输入规模，要记录环境，要把冷启动、预热、核心循环、提交和清理分开。如果工具权限不足，也要写清楚不能观察什么，而不是把猜测写成结论。系统教材的严谨性来自证据链：现象、假设、对照、指标、结论边界。

阻塞也有边界。过度拆池也会带来调度和资源配置复杂度。工程判断不是把一个技术推到所有地方，而是知道它在哪些约束下成立。读者在本章应形成一种习惯：每学到一个机制，都要问它解决了什么失败、引入了什么新成本、需要什么测试、在什么输入或部署环境下会失效。

\topic{上下文切换}

先把问题收窄到一个可推导的场景：为什么切换线程不只是保存寄存器。在把日志作业拆成多个可执行任务里，这个问题不会以术语形式出现，而是以结果不稳定、吞吐不增长、延迟抖动或恢复困难的形式出现。如果一上来只背上下文切换的定义，读者很容易把它当成孤立知识点；更好的顺序是先看朴素方案为什么合理，再看它在真实约束下怎样失败。

朴素方案通常是：认为切换成本很小。这个方案的吸引力在于它贴近单线程直觉，代码短，状态少，测试也容易写。但是它隐含了几个没有说出口的假设：输入总是干净，资源近似无限，执行不会交错，失败不会发生，硬件成本和源码结构大体一致。一旦这些假设被打破，缓存、TLB、分支历史和运行队列都会受到影响。这时继续在原方案上加 if 或加锁，往往只会把真正的问题埋得更深。

上下文切换就是从这个失败里长出来的概念。它的核心模型可以这样理解：上下文切换是执行上下文和局部性的迁移。这个模型必须同时放在三个层次看。源码层要说明谁读谁写、谁拥有对象、哪个状态允许变化；运行时层要说明线程、队列、任务和 I/O 怎样排队；硬件或系统层要说明缓存、页、调度、文件系统或网络怎样参与成本。三层缺一层，解释都会变形：只有源码层会看不见隐藏成本，只有硬件层会丢失业务语义，只有运行时层又容易把正确性和性能混在一起。

观察方法不能停在“跑一下变快了”。这一点在本章尤其重要：用 perf stat 记录 context switches，配合 trace 看等待。实验要有 reference，要固定输入规模，要记录环境，要把冷启动、预热、核心循环、提交和清理分开。如果工具权限不足，也要写清楚不能观察什么，而不是把猜测写成结论。系统教材的严谨性来自证据链：现象、假设、对照、指标、结论边界。

上下文切换也有边界。少量切换正常，问题是无意义高频切换。工程判断不是把一个技术推到所有地方，而是知道它在哪些约束下成立。读者在本章应形成一种习惯：每学到一个机制，都要问它解决了什么失败、引入了什么新成本、需要什么测试、在什么输入或部署环境下会失效。

\topic{亲和性}

先把问题收窄到一个可推导的场景：绑定线程为什么可能更稳定，也可能更慢。在把日志作业拆成多个可执行任务里，这个问题不会以术语形式出现，而是以结果不稳定、吞吐不增长、延迟抖动或恢复困难的形式出现。如果一上来只背亲和性的定义，读者很容易把它当成孤立知识点；更好的顺序是先看朴素方案为什么合理，再看它在真实约束下怎样失败。

朴素方案通常是：让操作系统自由调度或固定绑定都当成绝对答案。这个方案的吸引力在于它贴近单线程直觉，代码短，状态少，测试也容易写。但是它隐含了几个没有说出口的假设：输入总是干净，资源近似无限，执行不会交错，失败不会发生，硬件成本和源码结构大体一致。一旦这些假设被打破，迁移会破坏缓存局部性，错误绑定会压住同一核心。这时继续在原方案上加 if 或加锁，往往只会把真正的问题埋得更深。

亲和性就是从这个失败里长出来的概念。它的核心模型可以这样理解：亲和性把 worker、CPU、cache 和 NUMA 位置联系起来。这个模型必须同时放在三个层次看。源码层要说明谁读谁写、谁拥有对象、哪个状态允许变化；运行时层要说明线程、队列、任务和 I/O 怎样排队；硬件或系统层要说明缓存、页、调度、文件系统或网络怎样参与成本。三层缺一层，解释都会变形：只有源码层会看不见隐藏成本，只有硬件层会丢失业务语义，只有运行时层又容易把正确性和性能混在一起。

观察方法不能停在“跑一下变快了”。这一点在本章尤其重要：比较绑定、未绑定、SMT sibling 和跨节点。实验要有 reference，要固定输入规模，要记录环境，要把冷启动、预热、核心循环、提交和清理分开。如果工具权限不足，也要写清楚不能观察什么，而不是把猜测写成结论。系统教材的严谨性来自证据链：现象、假设、对照、指标、结论边界。

亲和性也有边界。生产环境还要考虑容器配额和其他进程。工程判断不是把一个技术推到所有地方，而是知道它在哪些约束下成立。读者在本章应形成一种习惯：每学到一个机制，都要问它解决了什么失败、引入了什么新成本、需要什么测试、在什么输入或部署环境下会失效。

\topic{大小核}

先把问题收窄到一个可推导的场景：手机或移动设备上线程数增加为什么可能变慢。在把日志作业拆成多个可执行任务里，这个问题不会以术语形式出现，而是以结果不稳定、吞吐不增长、延迟抖动或恢复困难的形式出现。如果一上来只背大小核的定义，读者很容易把它当成孤立知识点；更好的顺序是先看朴素方案为什么合理，再看它在真实约束下怎样失败。

朴素方案通常是：把所有核心视为同构。这个方案的吸引力在于它贴近单线程直觉，代码短，状态少，测试也容易写。但是它隐含了几个没有说出口的假设：输入总是干净，资源近似无限，执行不会交错，失败不会发生，硬件成本和源码结构大体一致。一旦这些假设被打破，小核、大核、温度和频率限制让扩展曲线不规则。这时继续在原方案上加 if 或加锁，往往只会把真正的问题埋得更深。

大小核就是从这个失败里长出来的概念。它的核心模型可以这样理解：异构核心需要记录核心类型、频率和调度策略。这个模型必须同时放在三个层次看。源码层要说明谁读谁写、谁拥有对象、哪个状态允许变化；运行时层要说明线程、队列、任务和 I/O 怎样排队；硬件或系统层要说明缓存、页、调度、文件系统或网络怎样参与成本。三层缺一层，解释都会变形：只有源码层会看不见隐藏成本，只有硬件层会丢失业务语义，只有运行时层又容易把正确性和性能混在一起。

观察方法不能停在“跑一下变快了”。这一点在本章尤其重要：在 Termux 环境记录 CPU 列表和频率变化。实验要有 reference，要固定输入规模，要记录环境，要把冷启动、预热、核心循环、提交和清理分开。如果工具权限不足，也要写清楚不能观察什么，而不是把猜测写成结论。系统教材的严谨性来自证据链：现象、假设、对照、指标、结论边界。

大小核也有边界。大小核不是 NUMA，不能混用解释。工程判断不是把一个技术推到所有地方，而是知道它在哪些约束下成立。读者在本章应形成一种习惯：每学到一个机制，都要问它解决了什么失败、引入了什么新成本、需要什么测试、在什么输入或部署环境下会失效。

\topic{线程池关闭}

先把问题收窄到一个可推导的场景：作业取消时，正在等待队列的 worker 如何退出。在把日志作业拆成多个可执行任务里，这个问题不会以术语形式出现，而是以结果不稳定、吞吐不增长、延迟抖动或恢复困难的形式出现。如果一上来只背线程池关闭的定义，读者很容易把它当成孤立知识点；更好的顺序是先看朴素方案为什么合理，再看它在真实约束下怎样失败。

朴素方案通常是：设置一个 bool 后等待线程自然醒来。这个方案的吸引力在于它贴近单线程直觉，代码短，状态少，测试也容易写。但是它隐含了几个没有说出口的假设：输入总是干净，资源近似无限，执行不会交错，失败不会发生，硬件成本和源码结构大体一致。一旦这些假设被打破，等待线程可能永远睡眠，任务可能半完成。这时继续在原方案上加 if 或加锁，往往只会把真正的问题埋得更深。

线程池关闭就是从这个失败里长出来的概念。它的核心模型可以这样理解：关闭语义要唤醒等待者、拒绝新任务、处理未完成任务并传播异常。这个模型必须同时放在三个层次看。源码层要说明谁读谁写、谁拥有对象、哪个状态允许变化；运行时层要说明线程、队列、任务和 I/O 怎样排队；硬件或系统层要说明缓存、页、调度、文件系统或网络怎样参与成本。三层缺一层，解释都会变形：只有源码层会看不见隐藏成本，只有硬件层会丢失业务语义，只有运行时层又容易把正确性和性能混在一起。

观察方法不能停在“跑一下变快了”。这一点在本章尤其重要：用关闭期间 push、pop、cancel 的测试矩阵验证。实验要有 reference，要固定输入规模，要记录环境，要把冷启动、预热、核心循环、提交和清理分开。如果工具权限不足，也要写清楚不能观察什么，而不是把猜测写成结论。系统教材的严谨性来自证据链：现象、假设、对照、指标、结论边界。

线程池关闭也有边界。粗暴 detach 线程会丢失生命周期管理。工程判断不是把一个技术推到所有地方，而是知道它在哪些约束下成立。读者在本章应形成一种习惯：每学到一个机制，都要问它解决了什么失败、引入了什么新成本、需要什么测试、在什么输入或部署环境下会失效。

\topic{调度 trace}

先把问题收窄到一个可推导的场景：为什么平均吞吐看不出某个分片拖尾。在把日志作业拆成多个可执行任务里，这个问题不会以术语形式出现，而是以结果不稳定、吞吐不增长、延迟抖动或恢复困难的形式出现。如果一上来只背调度 trace的定义，读者很容易把它当成孤立知识点；更好的顺序是先看朴素方案为什么合理，再看它在真实约束下怎样失败。

朴素方案通常是：只记录作业总耗时。这个方案的吸引力在于它贴近单线程直觉，代码短，状态少，测试也容易写。但是它隐含了几个没有说出口的假设：输入总是干净，资源近似无限，执行不会交错，失败不会发生，硬件成本和源码结构大体一致。一旦这些假设被打破，长尾任务、迁移和等待被平均值掩盖。这时继续在原方案上加 if 或加锁，往往只会把真正的问题埋得更深。

调度 trace就是从这个失败里长出来的概念。它的核心模型可以这样理解：调度 trace 记录任务从提交到完成的每个阶段。这个模型必须同时放在三个层次看。源码层要说明谁读谁写、谁拥有对象、哪个状态允许变化；运行时层要说明线程、队列、任务和 I/O 怎样排队；硬件或系统层要说明缓存、页、调度、文件系统或网络怎样参与成本。三层缺一层，解释都会变形：只有源码层会看不见隐藏成本，只有硬件层会丢失业务语义，只有运行时层又容易把正确性和性能混在一起。

观察方法不能停在“跑一下变快了”。这一点在本章尤其重要：输出 task id、worker id、start、end、wait 和状态。实验要有 reference，要固定输入规模，要记录环境，要把冷启动、预热、核心循环、提交和清理分开。如果工具权限不足，也要写清楚不能观察什么，而不是把猜测写成结论。系统教材的严谨性来自证据链：现象、假设、对照、指标、结论边界。

调度 trace也有边界。trace 不能无限细，热路径要控制开销。工程判断不是把一个技术推到所有地方，而是知道它在哪些约束下成立。读者在本章应形成一种习惯：每学到一个机制，都要问它解决了什么失败、引入了什么新成本、需要什么测试、在什么输入或部署环境下会失效。

\topic{案例与实验串联}

案例“每任务一线程”用于把抽象机制落到一段可复现实验。场景是：把一千个小分片分别创建线程处理。这个场景要足够小，小到读者可以手工画出数据流、状态变化和成本路径；又要足够真实，能暴露教材正文讨论的失败模式。

第一版不要急着写最终答案，而应保留一个朴素版本：直接 thread 构造并 join。朴素版本的价值不是性能，而是语义清楚。它告诉我们结果应该是什么，也告诉我们优化前系统在哪些地方偷懒。

第二版再引入改进：固定线程池处理任务队列。改进必须说明成本迁移到哪里，而不是只说“更快”。例如减少共享写可能增加最终合并成本，批处理可能增加尾延迟，分片可能引入负载倾斜，异步可能让生命周期更难验证。

验证时重点看：比较线程创建时间、切换次数和吞吐。报告里至少写出输入规模、硬件或系统环境、编译选项、运行命令、关键指标和反例。若结果与预期不同，优先检查实验变量，而不是马上改代码。
案例“阻塞混入计算池”用于把抽象机制落到一段可复现实验。场景是：部分任务需要等待文件 I/O，部分任务纯 CPU。这个场景要足够小，小到读者可以手工画出数据流、状态变化和成本路径；又要足够真实，能暴露教材正文讨论的失败模式。

第一版不要急着写最终答案，而应保留一个朴素版本：全部放同一个 worker 池。朴素版本的价值不是性能，而是语义清楚。它告诉我们结果应该是什么，也告诉我们优化前系统在哪些地方偷懒。

第二版再引入改进：隔离 I/O 或使用异步完成通知。改进必须说明成本迁移到哪里，而不是只说“更快”。例如减少共享写可能增加最终合并成本，批处理可能增加尾延迟，分片可能引入负载倾斜，异步可能让生命周期更难验证。

验证时重点看：观察计算任务是否被阻塞任务拖慢。报告里至少写出输入规模、硬件或系统环境、编译选项、运行命令、关键指标和反例。若结果与预期不同，优先检查实验变量，而不是马上改代码。
案例“移动设备扩展曲线”用于把抽象机制落到一段可复现实验。场景是：在 Termux 上从一线程逐步增加到硬件线程数。这个场景要足够小，小到读者可以手工画出数据流、状态变化和成本路径；又要足够真实，能暴露教材正文讨论的失败模式。

第一版不要急着写最终答案，而应保留一个朴素版本：只看最大线程数结果。朴素版本的价值不是性能，而是语义清楚。它告诉我们结果应该是什么，也告诉我们优化前系统在哪些地方偷懒。

第二版再引入改进：记录每个线程数、频率和温度趋势。改进必须说明成本迁移到哪里，而不是只说“更快”。例如减少共享写可能增加最终合并成本，批处理可能增加尾延迟，分片可能引入负载倾斜，异步可能让生命周期更难验证。

验证时重点看：区分大小核、降频和同步瓶颈。报告里至少写出输入规模、硬件或系统环境、编译选项、运行命令、关键指标和反例。若结果与预期不同，优先检查实验变量，而不是马上改代码。

\topic{Linux 源码阅读与系统观察}

Linux 源码阅读要从用户态现象倒推入口。本章可围绕 \filepath{kernel/sched/core.c}、\filepath{kernel/sched/fair.c}、\filepath{kernel/futex/core.c} 做窄范围阅读。不要试图一次读完整子系统，先选择一个问题，例如一次阻塞为什么睡眠、一次缺页怎样建立映射、一次唤醒怎样进入 run queue、一次写回怎样变成持久化边界。

阅读报告应固定内核版本、阅读路径和实验命令。第一段写用户态最小程序，第二段写观察到的现象，第三段写源码中的关键结构和状态变化，第四段写结论边界。若当前设备不是对应内核，报告必须说明源码只是机制参考，而不是当前运行系统的直接证据。

源码阅读还要看数据结构，而不只是函数名。队列、树、链表、位图、引用计数、状态枚举、等待队列、页表、inode、bio、task 结构这些细节，决定了机制怎样落地。读者要练习把一个用户态概念翻译成内核对象：线程对应 task，等待对应 wait queue 或 futex 路径，文件缓存对应 page cache，内存策略对应 VMA 和页面分配。

\topic{贯穿项目检查点}

把本章落到贯穿项目时，不要只加一个接口或一个 benchmark。每次新增能力都应有语义目标、最小实现、对照实验、指标和失败注入。项目不是堆功能，而是把本章的机制变成可运行、可检查、可复盘的工程对象。

\begin{lstlisting}[numbers=none]
chapter project checkpoints:
  1. 实现固定大小线程池
  2. 记录 worker id 和 task id
  3. 加入关闭和取消测试
  4. 比较每任务一线程和线程池
  5. 写亲和性与环境记录
\end{lstlisting}

项目报告要回答四个问题。第一，朴素版本是什么，为什么不足。第二，改进版本改变了哪个边界，是数据边界、执行边界、同步边界、I/O 边界还是提交边界。第三，证据是什么，哪些指标支持结论。第四，剩余风险是什么，在哪些输入、机器或失败条件下结论可能不成立。只有这样，项目才不会变成代码展示，而会变成系统能力训练。



\topic{第二轮深水区：把机制讲成可验证能力}

本章继续扩写时，重点不再是补充更多名词，而是把已经出现的机制变成可验证能力。围绕把日志作业拆成多个可执行任务，读者最终应能做三件事：解释为什么朴素方案失败，设计能隔离变量的实验，把实验结论落到代码边界。这三件事缺一不可。只会解释会变成空谈，只会跑实验会变成工具使用，只会改代码又会在复杂系统里丢掉语义。
本章的主失败是：线程数量、任务粒度、阻塞和调度位置没有边界，导致过度并行、上下文切换和不可复现实验。这个失败会在不同尺度反复出现。小输入里它可能只是一次结果不稳定；多线程里它可能变成锁竞争、乱序或重复写；I/O 和分布式里它可能变成提交不清、恢复不可靠或重试放大。所以本章的每个机制都要写出尺度变化后的表现。

\topic{深挖：任务和线程}

围绕“任务和线程”，先把它放回一个具体问题：一个输入分片为什么不应该必然对应一个系统线程。如果这个问题只在本机分布式模拟里出现，读者很容易把它当成局部技巧；但在把日志作业拆成多个可执行任务中，它会沿着输入、执行、同步、I/O 和提交一路传播。所以讲解时不能只写定义，而要让读者看到它怎样从一个小失败变成系统性约束。

朴素写法通常是“每个分片创建一个 thread”。这句话看似简单，却包含几个默认前提：状态只有一个拥有者，执行顺序可预测，资源足够近，失败不会穿过边界，观察结果能够代表真实成本。这些前提在教学小程序中可能成立，在把日志作业拆成多个可执行任务里却会被输入规模、线程交错、硬件拓扑或故障注入逐个打破。

失败信号是“线程创建和调度开销随分片数膨胀，错误传播和关闭也难控制”。注意，这个失败不一定表现为崩溃。它也可能表现为吞吐不再增长、p99 抖动、重启后结果多了一份、某个分区长尾、CPU 使用率很高但有效工作很少。教材要把这些信号和机制绑定起来，否则读者只会背术语，遇到真实现象时仍然不知道从哪里下手。

更可靠的推导顺序是先写出不变量。对任务和线程来说，不变量不是一句口号，而是本章要保护的事实：哪些数据只能写一次，哪些状态可以重试，哪些结果可以被缓存，哪些成本必须摊销，哪些错误必须外显。只要不变量没有写清，后面的代码、实验和优化都没有稳定坐标。

然后才引入模型：任务是工作单元，线程是执行资源，运行时负责映射。这个模型应同时解释正确性和成本。正确性部分回答“为什么结果可信”，成本部分回答“为什么这个实现会快或慢”。如果一个模型只能解释速度，却解释不了失败恢复，它不适合系统教材；如果只能解释语义，却完全不关心 cache、调度、I/O 或网络成本，也不适合第二册。

实验应围绕“只改变一个变量”设计。针对任务和线程，最小实验可以保留朴素版本，再实现一个改进版本，输入规模从小到大，线程数或并发度从一到目标机器上限，错误注入从无到有。每一组实验都应记录 reference 结果、核心指标、环境和命令。这样读者看到差异时，可以判断差异来自机制本身，而不是来自初始化、缓存冷热、调度迁移或文件系统状态。

观察手段是：比较每任务一线程和固定线程池。观察不是为了堆工具名，而是为了回答一个具体假设。若假设是共享写导致扩展性下降，就应看线程数曲线、写入布局和 cache 相关指标；若假设是提交边界错误，就应做崩溃点矩阵；若假设是队列背压，就应看水位、等待和上游速率。工具输出必须回到假设，不要把截图或命令输出当成结论。

最后要讲边界：任务太大负载不均，太小调度成本高。高质量教材必须把反例放在正文里。读者需要知道什么时候该用这个机制，什么时候它会制造新问题，什么时候应该退回更简单方案。比如一个单线程工具没有并发共享，强行引入复杂运行时只会增加维护成本；一个没有恢复要求的临时分析脚本，也不必承担完整 manifest 协议。

把这一点落实到代码审查，可以要求每个涉及任务和线程的改动都回答五个问题：朴素版本是什么，新增机制保护了哪个不变量，新增成本在哪里，如何回退，哪些测试证明边界。这五个问题比“用了某某技术”更能判断质量，也能防止团队在没有证据时追逐复杂实现。

\topic{深挖：过度订阅}

围绕“过度订阅”，先把它放回一个具体问题：线程数超过可用执行资源会发生什么。如果这个问题只在真实线上服务里出现，读者很容易把它当成局部技巧；但在把日志作业拆成多个可执行任务中，它会沿着输入、执行、同步、I/O 和提交一路传播。所以讲解时不能只写定义，而要让读者看到它怎样从一个小失败变成系统性约束。

朴素写法通常是“按输入分片数量创建线程”。这句话看似简单，却包含几个默认前提：状态只有一个拥有者，执行顺序可预测，资源足够近，失败不会穿过边界，观察结果能够代表真实成本。这些前提在教学小程序中可能成立，在把日志作业拆成多个可执行任务里却会被输入规模、线程交错、硬件拓扑或故障注入逐个打破。

失败信号是“上下文切换、缓存失效和调度延迟增加”。注意，这个失败不一定表现为崩溃。它也可能表现为吞吐不再增长、p99 抖动、重启后结果多了一份、某个分区长尾、CPU 使用率很高但有效工作很少。教材要把这些信号和机制绑定起来，否则读者只会背术语，遇到真实现象时仍然不知道从哪里下手。

更可靠的推导顺序是先写出不变量。对过度订阅来说，不变量不是一句口号，而是本章要保护的事实：哪些数据只能写一次，哪些状态可以重试，哪些结果可以被缓存，哪些成本必须摊销，哪些错误必须外显。只要不变量没有写清，后面的代码、实验和优化都没有稳定坐标。

然后才引入模型：过度订阅让 runnable 线程争夺有限 CPU。这个模型应同时解释正确性和成本。正确性部分回答“为什么结果可信”，成本部分回答“为什么这个实现会快或慢”。如果一个模型只能解释速度，却解释不了失败恢复，它不适合系统教材；如果只能解释语义，却完全不关心 cache、调度、I/O 或网络成本，也不适合第二册。

实验应围绕“只改变一个变量”设计。针对过度订阅，最小实验可以保留朴素版本，再实现一个改进版本，输入规模从小到大，线程数或并发度从一到目标机器上限，错误注入从无到有。每一组实验都应记录 reference 结果、核心指标、环境和命令。这样读者看到差异时，可以判断差异来自机制本身，而不是来自初始化、缓存冷热、调度迁移或文件系统状态。

观察手段是：记录 runnable 数、context switch 和吞吐曲线。观察不是为了堆工具名，而是为了回答一个具体假设。若假设是共享写导致扩展性下降，就应看线程数曲线、写入布局和 cache 相关指标；若假设是提交边界错误，就应做崩溃点矩阵；若假设是队列背压，就应看水位、等待和上游速率。工具输出必须回到假设，不要把截图或命令输出当成结论。

最后要讲边界：I/O 阻塞任务可能需要不同策略，不能只按核心数固定。高质量教材必须把反例放在正文里。读者需要知道什么时候该用这个机制，什么时候它会制造新问题，什么时候应该退回更简单方案。比如一个单线程工具没有并发共享，强行引入复杂运行时只会增加维护成本；一个没有恢复要求的临时分析脚本，也不必承担完整 manifest 协议。

把这一点落实到代码审查，可以要求每个涉及过度订阅的改动都回答五个问题：朴素版本是什么，新增机制保护了哪个不变量，新增成本在哪里，如何回退，哪些测试证明边界。这五个问题比“用了某某技术”更能判断质量，也能防止团队在没有证据时追逐复杂实现。

\topic{深挖：阻塞}

围绕“阻塞”，先把它放回一个具体问题：一个 worker 阻塞在 I/O 或条件变量上是否还算占用计算资源。如果这个问题只在单线程小输入里出现，读者很容易把它当成局部技巧；但在把日志作业拆成多个可执行任务中，它会沿着输入、执行、同步、I/O 和提交一路传播。所以讲解时不能只写定义，而要让读者看到它怎样从一个小失败变成系统性约束。

朴素写法通常是“把阻塞任务和计算任务放同一个池”。这句话看似简单，却包含几个默认前提：状态只有一个拥有者，执行顺序可预测，资源足够近，失败不会穿过边界，观察结果能够代表真实成本。这些前提在教学小程序中可能成立，在把日志作业拆成多个可执行任务里却会被输入规模、线程交错、硬件拓扑或故障注入逐个打破。

失败信号是“计算线程被等待占住，吞吐下降且排队难解释”。注意，这个失败不一定表现为崩溃。它也可能表现为吞吐不再增长、p99 抖动、重启后结果多了一份、某个分区长尾、CPU 使用率很高但有效工作很少。教材要把这些信号和机制绑定起来，否则读者只会背术语，遇到真实现象时仍然不知道从哪里下手。

更可靠的推导顺序是先写出不变量。对阻塞来说，不变量不是一句口号，而是本章要保护的事实：哪些数据只能写一次，哪些状态可以重试，哪些结果可以被缓存，哪些成本必须摊销，哪些错误必须外显。只要不变量没有写清，后面的代码、实验和优化都没有稳定坐标。

然后才引入模型：阻塞会改变线程池容量模型，需要隔离或异步化。这个模型应同时解释正确性和成本。正确性部分回答“为什么结果可信”，成本部分回答“为什么这个实现会快或慢”。如果一个模型只能解释速度，却解释不了失败恢复，它不适合系统教材；如果只能解释语义，却完全不关心 cache、调度、I/O 或网络成本，也不适合第二册。

实验应围绕“只改变一个变量”设计。针对阻塞，最小实验可以保留朴素版本，再实现一个改进版本，输入规模从小到大，线程数或并发度从一到目标机器上限，错误注入从无到有。每一组实验都应记录 reference 结果、核心指标、环境和命令。这样读者看到差异时，可以判断差异来自机制本身，而不是来自初始化、缓存冷热、调度迁移或文件系统状态。

观察手段是：记录线程状态、队列长度和等待原因。观察不是为了堆工具名，而是为了回答一个具体假设。若假设是共享写导致扩展性下降，就应看线程数曲线、写入布局和 cache 相关指标；若假设是提交边界错误，就应做崩溃点矩阵；若假设是队列背压，就应看水位、等待和上游速率。工具输出必须回到假设，不要把截图或命令输出当成结论。

最后要讲边界：过度拆池也会带来调度和资源配置复杂度。高质量教材必须把反例放在正文里。读者需要知道什么时候该用这个机制，什么时候它会制造新问题，什么时候应该退回更简单方案。比如一个单线程工具没有并发共享，强行引入复杂运行时只会增加维护成本；一个没有恢复要求的临时分析脚本，也不必承担完整 manifest 协议。

把这一点落实到代码审查，可以要求每个涉及阻塞的改动都回答五个问题：朴素版本是什么，新增机制保护了哪个不变量，新增成本在哪里，如何回退，哪些测试证明边界。这五个问题比“用了某某技术”更能判断质量，也能防止团队在没有证据时追逐复杂实现。

\topic{深挖：上下文切换}

围绕“上下文切换”，先把它放回一个具体问题：为什么切换线程不只是保存寄存器。如果这个问题只在单机多线程里出现，读者很容易把它当成局部技巧；但在把日志作业拆成多个可执行任务中，它会沿着输入、执行、同步、I/O 和提交一路传播。所以讲解时不能只写定义，而要让读者看到它怎样从一个小失败变成系统性约束。

朴素写法通常是“认为切换成本很小”。这句话看似简单，却包含几个默认前提：状态只有一个拥有者，执行顺序可预测，资源足够近，失败不会穿过边界，观察结果能够代表真实成本。这些前提在教学小程序中可能成立，在把日志作业拆成多个可执行任务里却会被输入规模、线程交错、硬件拓扑或故障注入逐个打破。

失败信号是“缓存、TLB、分支历史和运行队列都会受到影响”。注意，这个失败不一定表现为崩溃。它也可能表现为吞吐不再增长、p99 抖动、重启后结果多了一份、某个分区长尾、CPU 使用率很高但有效工作很少。教材要把这些信号和机制绑定起来，否则读者只会背术语，遇到真实现象时仍然不知道从哪里下手。

更可靠的推导顺序是先写出不变量。对上下文切换来说，不变量不是一句口号，而是本章要保护的事实：哪些数据只能写一次，哪些状态可以重试，哪些结果可以被缓存，哪些成本必须摊销，哪些错误必须外显。只要不变量没有写清，后面的代码、实验和优化都没有稳定坐标。

然后才引入模型：上下文切换是执行上下文和局部性的迁移。这个模型应同时解释正确性和成本。正确性部分回答“为什么结果可信”，成本部分回答“为什么这个实现会快或慢”。如果一个模型只能解释速度，却解释不了失败恢复，它不适合系统教材；如果只能解释语义，却完全不关心 cache、调度、I/O 或网络成本，也不适合第二册。

实验应围绕“只改变一个变量”设计。针对上下文切换，最小实验可以保留朴素版本，再实现一个改进版本，输入规模从小到大，线程数或并发度从一到目标机器上限，错误注入从无到有。每一组实验都应记录 reference 结果、核心指标、环境和命令。这样读者看到差异时，可以判断差异来自机制本身，而不是来自初始化、缓存冷热、调度迁移或文件系统状态。

观察手段是：用 perf stat 记录 context switches，配合 trace 看等待。观察不是为了堆工具名，而是为了回答一个具体假设。若假设是共享写导致扩展性下降，就应看线程数曲线、写入布局和 cache 相关指标；若假设是提交边界错误，就应做崩溃点矩阵；若假设是队列背压，就应看水位、等待和上游速率。工具输出必须回到假设，不要把截图或命令输出当成结论。

最后要讲边界：少量切换正常，问题是无意义高频切换。高质量教材必须把反例放在正文里。读者需要知道什么时候该用这个机制，什么时候它会制造新问题，什么时候应该退回更简单方案。比如一个单线程工具没有并发共享，强行引入复杂运行时只会增加维护成本；一个没有恢复要求的临时分析脚本，也不必承担完整 manifest 协议。

把这一点落实到代码审查，可以要求每个涉及上下文切换的改动都回答五个问题：朴素版本是什么，新增机制保护了哪个不变量，新增成本在哪里，如何回退，哪些测试证明边界。这五个问题比“用了某某技术”更能判断质量，也能防止团队在没有证据时追逐复杂实现。

\topic{深挖：亲和性}

围绕“亲和性”，先把它放回一个具体问题：绑定线程为什么可能更稳定，也可能更慢。如果这个问题只在NUMA 或异构核心里出现，读者很容易把它当成局部技巧；但在把日志作业拆成多个可执行任务中，它会沿着输入、执行、同步、I/O 和提交一路传播。所以讲解时不能只写定义，而要让读者看到它怎样从一个小失败变成系统性约束。

朴素写法通常是“让操作系统自由调度或固定绑定都当成绝对答案”。这句话看似简单，却包含几个默认前提：状态只有一个拥有者，执行顺序可预测，资源足够近，失败不会穿过边界，观察结果能够代表真实成本。这些前提在教学小程序中可能成立，在把日志作业拆成多个可执行任务里却会被输入规模、线程交错、硬件拓扑或故障注入逐个打破。

失败信号是“迁移会破坏缓存局部性，错误绑定会压住同一核心”。注意，这个失败不一定表现为崩溃。它也可能表现为吞吐不再增长、p99 抖动、重启后结果多了一份、某个分区长尾、CPU 使用率很高但有效工作很少。教材要把这些信号和机制绑定起来，否则读者只会背术语，遇到真实现象时仍然不知道从哪里下手。

更可靠的推导顺序是先写出不变量。对亲和性来说，不变量不是一句口号，而是本章要保护的事实：哪些数据只能写一次，哪些状态可以重试，哪些结果可以被缓存，哪些成本必须摊销，哪些错误必须外显。只要不变量没有写清，后面的代码、实验和优化都没有稳定坐标。

然后才引入模型：亲和性把 worker、CPU、cache 和 NUMA 位置联系起来。这个模型应同时解释正确性和成本。正确性部分回答“为什么结果可信”，成本部分回答“为什么这个实现会快或慢”。如果一个模型只能解释速度，却解释不了失败恢复，它不适合系统教材；如果只能解释语义，却完全不关心 cache、调度、I/O 或网络成本，也不适合第二册。

实验应围绕“只改变一个变量”设计。针对亲和性，最小实验可以保留朴素版本，再实现一个改进版本，输入规模从小到大，线程数或并发度从一到目标机器上限，错误注入从无到有。每一组实验都应记录 reference 结果、核心指标、环境和命令。这样读者看到差异时，可以判断差异来自机制本身，而不是来自初始化、缓存冷热、调度迁移或文件系统状态。

观察手段是：比较绑定、未绑定、SMT sibling 和跨节点。观察不是为了堆工具名，而是为了回答一个具体假设。若假设是共享写导致扩展性下降，就应看线程数曲线、写入布局和 cache 相关指标；若假设是提交边界错误，就应做崩溃点矩阵；若假设是队列背压，就应看水位、等待和上游速率。工具输出必须回到假设，不要把截图或命令输出当成结论。

最后要讲边界：生产环境还要考虑容器配额和其他进程。高质量教材必须把反例放在正文里。读者需要知道什么时候该用这个机制，什么时候它会制造新问题，什么时候应该退回更简单方案。比如一个单线程工具没有并发共享，强行引入复杂运行时只会增加维护成本；一个没有恢复要求的临时分析脚本，也不必承担完整 manifest 协议。

把这一点落实到代码审查，可以要求每个涉及亲和性的改动都回答五个问题：朴素版本是什么，新增机制保护了哪个不变量，新增成本在哪里，如何回退，哪些测试证明边界。这五个问题比“用了某某技术”更能判断质量，也能防止团队在没有证据时追逐复杂实现。

\topic{深挖：大小核}

围绕“大小核”，先把它放回一个具体问题：手机或移动设备上线程数增加为什么可能变慢。如果这个问题只在本机分布式模拟里出现，读者很容易把它当成局部技巧；但在把日志作业拆成多个可执行任务中，它会沿着输入、执行、同步、I/O 和提交一路传播。所以讲解时不能只写定义，而要让读者看到它怎样从一个小失败变成系统性约束。

朴素写法通常是“把所有核心视为同构”。这句话看似简单，却包含几个默认前提：状态只有一个拥有者，执行顺序可预测，资源足够近，失败不会穿过边界，观察结果能够代表真实成本。这些前提在教学小程序中可能成立，在把日志作业拆成多个可执行任务里却会被输入规模、线程交错、硬件拓扑或故障注入逐个打破。

失败信号是“小核、大核、温度和频率限制让扩展曲线不规则”。注意，这个失败不一定表现为崩溃。它也可能表现为吞吐不再增长、p99 抖动、重启后结果多了一份、某个分区长尾、CPU 使用率很高但有效工作很少。教材要把这些信号和机制绑定起来，否则读者只会背术语，遇到真实现象时仍然不知道从哪里下手。

更可靠的推导顺序是先写出不变量。对大小核来说，不变量不是一句口号，而是本章要保护的事实：哪些数据只能写一次，哪些状态可以重试，哪些结果可以被缓存，哪些成本必须摊销，哪些错误必须外显。只要不变量没有写清，后面的代码、实验和优化都没有稳定坐标。

然后才引入模型：异构核心需要记录核心类型、频率和调度策略。这个模型应同时解释正确性和成本。正确性部分回答“为什么结果可信”，成本部分回答“为什么这个实现会快或慢”。如果一个模型只能解释速度，却解释不了失败恢复，它不适合系统教材；如果只能解释语义，却完全不关心 cache、调度、I/O 或网络成本，也不适合第二册。

实验应围绕“只改变一个变量”设计。针对大小核，最小实验可以保留朴素版本，再实现一个改进版本，输入规模从小到大，线程数或并发度从一到目标机器上限，错误注入从无到有。每一组实验都应记录 reference 结果、核心指标、环境和命令。这样读者看到差异时，可以判断差异来自机制本身，而不是来自初始化、缓存冷热、调度迁移或文件系统状态。

观察手段是：在 Termux 环境记录 CPU 列表和频率变化。观察不是为了堆工具名，而是为了回答一个具体假设。若假设是共享写导致扩展性下降，就应看线程数曲线、写入布局和 cache 相关指标；若假设是提交边界错误，就应做崩溃点矩阵；若假设是队列背压，就应看水位、等待和上游速率。工具输出必须回到假设，不要把截图或命令输出当成结论。

最后要讲边界：大小核不是 NUMA，不能混用解释。高质量教材必须把反例放在正文里。读者需要知道什么时候该用这个机制，什么时候它会制造新问题，什么时候应该退回更简单方案。比如一个单线程工具没有并发共享，强行引入复杂运行时只会增加维护成本；一个没有恢复要求的临时分析脚本，也不必承担完整 manifest 协议。

把这一点落实到代码审查，可以要求每个涉及大小核的改动都回答五个问题：朴素版本是什么，新增机制保护了哪个不变量，新增成本在哪里，如何回退，哪些测试证明边界。这五个问题比“用了某某技术”更能判断质量，也能防止团队在没有证据时追逐复杂实现。

\topic{深挖：线程池关闭}

围绕“线程池关闭”，先把它放回一个具体问题：作业取消时，正在等待队列的 worker 如何退出。如果这个问题只在真实线上服务里出现，读者很容易把它当成局部技巧；但在把日志作业拆成多个可执行任务中，它会沿着输入、执行、同步、I/O 和提交一路传播。所以讲解时不能只写定义，而要让读者看到它怎样从一个小失败变成系统性约束。

朴素写法通常是“设置一个 bool 后等待线程自然醒来”。这句话看似简单，却包含几个默认前提：状态只有一个拥有者，执行顺序可预测，资源足够近，失败不会穿过边界，观察结果能够代表真实成本。这些前提在教学小程序中可能成立，在把日志作业拆成多个可执行任务里却会被输入规模、线程交错、硬件拓扑或故障注入逐个打破。

失败信号是“等待线程可能永远睡眠，任务可能半完成”。注意，这个失败不一定表现为崩溃。它也可能表现为吞吐不再增长、p99 抖动、重启后结果多了一份、某个分区长尾、CPU 使用率很高但有效工作很少。教材要把这些信号和机制绑定起来，否则读者只会背术语，遇到真实现象时仍然不知道从哪里下手。

更可靠的推导顺序是先写出不变量。对线程池关闭来说，不变量不是一句口号，而是本章要保护的事实：哪些数据只能写一次，哪些状态可以重试，哪些结果可以被缓存，哪些成本必须摊销，哪些错误必须外显。只要不变量没有写清，后面的代码、实验和优化都没有稳定坐标。

然后才引入模型：关闭语义要唤醒等待者、拒绝新任务、处理未完成任务并传播异常。这个模型应同时解释正确性和成本。正确性部分回答“为什么结果可信”，成本部分回答“为什么这个实现会快或慢”。如果一个模型只能解释速度，却解释不了失败恢复，它不适合系统教材；如果只能解释语义，却完全不关心 cache、调度、I/O 或网络成本，也不适合第二册。

实验应围绕“只改变一个变量”设计。针对线程池关闭，最小实验可以保留朴素版本，再实现一个改进版本，输入规模从小到大，线程数或并发度从一到目标机器上限，错误注入从无到有。每一组实验都应记录 reference 结果、核心指标、环境和命令。这样读者看到差异时，可以判断差异来自机制本身，而不是来自初始化、缓存冷热、调度迁移或文件系统状态。

观察手段是：用关闭期间 push、pop、cancel 的测试矩阵验证。观察不是为了堆工具名，而是为了回答一个具体假设。若假设是共享写导致扩展性下降，就应看线程数曲线、写入布局和 cache 相关指标；若假设是提交边界错误，就应做崩溃点矩阵；若假设是队列背压，就应看水位、等待和上游速率。工具输出必须回到假设，不要把截图或命令输出当成结论。

最后要讲边界：粗暴 detach 线程会丢失生命周期管理。高质量教材必须把反例放在正文里。读者需要知道什么时候该用这个机制，什么时候它会制造新问题，什么时候应该退回更简单方案。比如一个单线程工具没有并发共享，强行引入复杂运行时只会增加维护成本；一个没有恢复要求的临时分析脚本，也不必承担完整 manifest 协议。

把这一点落实到代码审查，可以要求每个涉及线程池关闭的改动都回答五个问题：朴素版本是什么，新增机制保护了哪个不变量，新增成本在哪里，如何回退，哪些测试证明边界。这五个问题比“用了某某技术”更能判断质量，也能防止团队在没有证据时追逐复杂实现。

\topic{深挖：调度 trace}

围绕“调度 trace”，先把它放回一个具体问题：为什么平均吞吐看不出某个分片拖尾。如果这个问题只在单线程小输入里出现，读者很容易把它当成局部技巧；但在把日志作业拆成多个可执行任务中，它会沿着输入、执行、同步、I/O 和提交一路传播。所以讲解时不能只写定义，而要让读者看到它怎样从一个小失败变成系统性约束。

朴素写法通常是“只记录作业总耗时”。这句话看似简单，却包含几个默认前提：状态只有一个拥有者，执行顺序可预测，资源足够近，失败不会穿过边界，观察结果能够代表真实成本。这些前提在教学小程序中可能成立，在把日志作业拆成多个可执行任务里却会被输入规模、线程交错、硬件拓扑或故障注入逐个打破。

失败信号是“长尾任务、迁移和等待被平均值掩盖”。注意，这个失败不一定表现为崩溃。它也可能表现为吞吐不再增长、p99 抖动、重启后结果多了一份、某个分区长尾、CPU 使用率很高但有效工作很少。教材要把这些信号和机制绑定起来，否则读者只会背术语，遇到真实现象时仍然不知道从哪里下手。

更可靠的推导顺序是先写出不变量。对调度 trace来说，不变量不是一句口号，而是本章要保护的事实：哪些数据只能写一次，哪些状态可以重试，哪些结果可以被缓存，哪些成本必须摊销，哪些错误必须外显。只要不变量没有写清，后面的代码、实验和优化都没有稳定坐标。

然后才引入模型：调度 trace 记录任务从提交到完成的每个阶段。这个模型应同时解释正确性和成本。正确性部分回答“为什么结果可信”，成本部分回答“为什么这个实现会快或慢”。如果一个模型只能解释速度，却解释不了失败恢复，它不适合系统教材；如果只能解释语义，却完全不关心 cache、调度、I/O 或网络成本，也不适合第二册。

实验应围绕“只改变一个变量”设计。针对调度 trace，最小实验可以保留朴素版本，再实现一个改进版本，输入规模从小到大，线程数或并发度从一到目标机器上限，错误注入从无到有。每一组实验都应记录 reference 结果、核心指标、环境和命令。这样读者看到差异时，可以判断差异来自机制本身，而不是来自初始化、缓存冷热、调度迁移或文件系统状态。

观察手段是：输出 task id、worker id、start、end、wait 和状态。观察不是为了堆工具名，而是为了回答一个具体假设。若假设是共享写导致扩展性下降，就应看线程数曲线、写入布局和 cache 相关指标；若假设是提交边界错误，就应做崩溃点矩阵；若假设是队列背压，就应看水位、等待和上游速率。工具输出必须回到假设，不要把截图或命令输出当成结论。

最后要讲边界：trace 不能无限细，热路径要控制开销。高质量教材必须把反例放在正文里。读者需要知道什么时候该用这个机制，什么时候它会制造新问题，什么时候应该退回更简单方案。比如一个单线程工具没有并发共享，强行引入复杂运行时只会增加维护成本；一个没有恢复要求的临时分析脚本，也不必承担完整 manifest 协议。

把这一点落实到代码审查，可以要求每个涉及调度 trace的改动都回答五个问题：朴素版本是什么，新增机制保护了哪个不变量，新增成本在哪里，如何回退，哪些测试证明边界。这五个问题比“用了某某技术”更能判断质量，也能防止团队在没有证据时追逐复杂实现。

\topic{案例深化：从朴素版到工程版}

案例深化“每任务一线程”要按三次迭代写。第一次只写 reference：把一千个小分片分别创建线程处理。reference 的代码可以慢，但必须把输入、输出、错误和边界写清。第二次写朴素工程版本：直接 thread 构造并 join。这一步不能省，因为读者需要看到直觉方案为什么诱人，也需要有一个失败对象。

第三次才写改进版本：固定线程池处理任务队列。改进说明要避免“因为更高级所以更快”这种表达，而要讲清它改变了哪条路径：减少了数据移动、减少了共享写、减少了系统调用、减少了重试放大，还是把不可恢复状态变成可恢复状态。

验收方式是：比较线程创建时间、切换次数和吞吐。验收报告至少包含一张对照表，一列是朴素版本，一列是改进版本，一列是反例。反例很重要，因为它能告诉读者改进不是什么时候都正确。如果一个案例没有反例，往往说明分析还停在宣传层面。
案例深化“阻塞混入计算池”要按三次迭代写。第一次只写 reference：部分任务需要等待文件 I/O，部分任务纯 CPU。reference 的代码可以慢，但必须把输入、输出、错误和边界写清。第二次写朴素工程版本：全部放同一个 worker 池。这一步不能省，因为读者需要看到直觉方案为什么诱人，也需要有一个失败对象。

第三次才写改进版本：隔离 I/O 或使用异步完成通知。改进说明要避免“因为更高级所以更快”这种表达，而要讲清它改变了哪条路径：减少了数据移动、减少了共享写、减少了系统调用、减少了重试放大，还是把不可恢复状态变成可恢复状态。

验收方式是：观察计算任务是否被阻塞任务拖慢。验收报告至少包含一张对照表，一列是朴素版本，一列是改进版本，一列是反例。反例很重要，因为它能告诉读者改进不是什么时候都正确。如果一个案例没有反例，往往说明分析还停在宣传层面。
案例深化“移动设备扩展曲线”要按三次迭代写。第一次只写 reference：在 Termux 上从一线程逐步增加到硬件线程数。reference 的代码可以慢，但必须把输入、输出、错误和边界写清。第二次写朴素工程版本：只看最大线程数结果。这一步不能省，因为读者需要看到直觉方案为什么诱人，也需要有一个失败对象。

第三次才写改进版本：记录每个线程数、频率和温度趋势。改进说明要避免“因为更高级所以更快”这种表达，而要讲清它改变了哪条路径：减少了数据移动、减少了共享写、减少了系统调用、减少了重试放大，还是把不可恢复状态变成可恢复状态。

验收方式是：区分大小核、降频和同步瓶颈。验收报告至少包含一张对照表，一列是朴素版本，一列是改进版本，一列是反例。反例很重要，因为它能告诉读者改进不是什么时候都正确。如果一个案例没有反例，往往说明分析还停在宣传层面。

\topic{实验矩阵：让结论可复现}

围绕任务和线程，实验矩阵至少包含正常输入、边界输入、压力输入和错误输入四类。正常输入验证功能，边界输入验证不变量，压力输入验证成本模型，错误输入验证恢复和报告。其中压力输入不只是“大”，还要能触发本章关心的形状：随机访问、热点 key、短任务、慢 I/O、重复消息或跨节点访问。

围绕过度订阅，实验矩阵至少包含正常输入、边界输入、压力输入和错误输入四类。正常输入验证功能，边界输入验证不变量，压力输入验证成本模型，错误输入验证恢复和报告。其中压力输入不只是“大”，还要能触发本章关心的形状：随机访问、热点 key、短任务、慢 I/O、重复消息或跨节点访问。

围绕阻塞，实验矩阵至少包含正常输入、边界输入、压力输入和错误输入四类。正常输入验证功能，边界输入验证不变量，压力输入验证成本模型，错误输入验证恢复和报告。其中压力输入不只是“大”，还要能触发本章关心的形状：随机访问、热点 key、短任务、慢 I/O、重复消息或跨节点访问。

围绕上下文切换，实验矩阵至少包含正常输入、边界输入、压力输入和错误输入四类。正常输入验证功能，边界输入验证不变量，压力输入验证成本模型，错误输入验证恢复和报告。其中压力输入不只是“大”，还要能触发本章关心的形状：随机访问、热点 key、短任务、慢 I/O、重复消息或跨节点访问。

围绕亲和性，实验矩阵至少包含正常输入、边界输入、压力输入和错误输入四类。正常输入验证功能，边界输入验证不变量，压力输入验证成本模型，错误输入验证恢复和报告。其中压力输入不只是“大”，还要能触发本章关心的形状：随机访问、热点 key、短任务、慢 I/O、重复消息或跨节点访问。

围绕大小核，实验矩阵至少包含正常输入、边界输入、压力输入和错误输入四类。正常输入验证功能，边界输入验证不变量，压力输入验证成本模型，错误输入验证恢复和报告。其中压力输入不只是“大”，还要能触发本章关心的形状：随机访问、热点 key、短任务、慢 I/O、重复消息或跨节点访问。

围绕线程池关闭，实验矩阵至少包含正常输入、边界输入、压力输入和错误输入四类。正常输入验证功能，边界输入验证不变量，压力输入验证成本模型，错误输入验证恢复和报告。其中压力输入不只是“大”，还要能触发本章关心的形状：随机访问、热点 key、短任务、慢 I/O、重复消息或跨节点访问。

围绕调度 trace，实验矩阵至少包含正常输入、边界输入、压力输入和错误输入四类。正常输入验证功能，边界输入验证不变量，压力输入验证成本模型，错误输入验证恢复和报告。其中压力输入不只是“大”，还要能触发本章关心的形状：随机访问、热点 key、短任务、慢 I/O、重复消息或跨节点访问。

\topic{Linux 源码阅读深化}

本章的 Linux 阅读入口仍然保持窄范围：\filepath{kernel/sched/core.c}、\filepath{kernel/sched/fair.c}、\filepath{kernel/futex/core.c}。阅读时不要追求覆盖率，而要追求因果链。从一个用户态现象出发，找到内核中的状态对象，再看状态怎样被修改、等待怎样被挂起、唤醒怎样发生、错误怎样返回。

源码阅读可以按四栏笔记整理。第一栏写用户态操作，例如线程等待、缺页、写文件、发送消息或计时。第二栏写内核对象，例如 task、VMA、page、inode、wait queue、bio、socket buffer。第三栏写关键状态变化，例如 runnable 到 sleeping、clean page 到 dirty page、pending task 到 committed task。第四栏写可观察指标或实验，例如 context switch、page fault、writeback、queue depth、retry count。

不要把 Linux 源码当成术语来源。源码阅读的目标是训练映射能力：把书中的抽象边界映射到真实系统里的结构和状态。读者不需要一次读懂整个内核，但要能解释一个最小现象为什么发生、哪些路径参与、哪些结论受当前内核版本和硬件限制。

\topic{项目验收：把本章能力写进仓库}

贯穿项目的验收不能只说“功能跑通”。本章至少要留下一个可重复实验、一个失败注入、一个指标输出和一个设计说明。实验说明机制是否真的被触发；失败注入说明边界是否真的被处理；指标输出说明结论是否有证据；设计说明让后续章节能复用这个模块。

本章项目检查点可以具体化为：1. 实现固定大小线程池
2. 记录 worker id 和 task id
3. 加入关闭和取消测试
4. 比较每任务一线程和线程池
5. 写亲和性与环境记录。这些检查点要进入仓库，而不是停留在阅读笔记。如果某个检查点因为当前设备权限、硬件或系统 API 不支持而无法完成，也要写清降级方案和未验证风险。

项目验收还应要求读者写一段复盘：本章新增机制让系统多了什么能力，也让系统多了什么复杂度。比如加入背压后内存更稳定，但生产者会等待；加入 route epoch 后旧请求能被拒绝，但所有消息都必须携带版本；加入 SIMD 后热内核更快，但尾部和数值语义要额外测试。这种复盘能把知识从“会用”推进到“会判断”。



\topic{最终收束：本章应形成的判断力}

读完本章，读者不应只记住任务和线程、调度 trace这些词，而应能围绕把日志作业拆成多个可执行任务做一次完整判断。完整判断从问题开始：现象是什么，朴素方案是什么，失败信号是什么，保护的不变量是什么，成本从哪里移动到哪里。若无法回答这些问题，就说明还停留在概念层，没有真正进入系统层。

本章最重要的反例是：线程数量、任务粒度、阻塞和调度位置没有边界，导致过度并行、上下文切换和不可复现实验。遇到类似反例时，第一步不是换技术，而是缩小问题。先固定输入，固定环境，保留 reference，再一次只改变一个变量。如果改动同时改变布局、线程数、批大小、同步方式和提交策略，即使结果变快，也无法知道原因。高质量工程实验必须克制变量。

以案例“每任务一线程”为例，最终报告应包含三段。第一段写把一千个小分片分别创建线程处理，说明读者要解决的不是抽象题，而是一个有输入、有状态、有失败方式的任务。第二段写朴素版本：直接 thread 构造并 join，并诚实说明它为什么吸引人。第三段写改进版本：固定线程池处理任务队列，再用比较线程创建时间、切换次数和吞吐验收。报告中若没有朴素版本，读者会误以为最终设计是凭空出现；若没有反例，读者会误以为最终设计永远正确。

本章还应留下一个审查表。正确性方面，检查输入合同、状态转移、所有权、重复执行、关闭和恢复。性能方面，检查数据移动、共享写、排队、系统调用、批大小、缓存冷热和拓扑。可观测性方面，检查是否有阶段耗时、队列水位、错误分类、任务身份和原始样本。每一项都要能落到代码、测试或报告，不要只停在口头承诺。

贯穿项目里，本章至少要完成这些检查点：实现固定大小线程池、记录 worker id 和 task id、加入关闭和取消测试、比较每任务一线程和线程池、写亲和性与环境记录。完成后不要急着进入下一章，要先写一页复盘：本章新增能力解决了什么失败，新增了哪些复杂度，哪些结论只在当前机器上验证，哪些结论可以迁移到更大系统。这种复盘会让整本书的知识保持连续，不会变成每章讲完就散掉的材料。

最后，用一句话收束本章：系统能力来自清楚的边界、可验证的不变量和可复现的证据。无论本章讨论的是硬件执行、内存层级、同步、运行时、I/O 还是分布式协议，只要缺少边界、不变量和证据，复杂实现就会变成风险来源。反过来，只要这三件事稳住，读者就能把一个朴素程序逐步推进成可靠的计算系统。




\topic{过线补充：常见误区、验收题和迁移能力}

本章最容易出现的误区，是把任务和线程、过度订阅、阻塞、上下文切换当成互相独立的知识点。在把日志作业拆成多个可执行任务中，它们其实围绕同一个失败展开：线程数量、任务粒度、阻塞和调度位置没有边界，导致过度并行、上下文切换和不可复现实验。如果读者只能分别解释这些概念，却不能说明它们在同一条数据流里怎样互相影响，就还没有真正掌握本章。例如一个优化减少了计算时间，却增加了队列等待；一个同步方案修复了正确性，却制造了共享写热点；一个提交协议提高了可靠性，却增加了持久化延迟。这些都不是矛盾，而是系统设计中的成本迁移。

本章的验收题可以这样设计：先给出一个最小版本的把日志作业拆成多个可执行任务，要求读者写出 reference；再引入一个压力条件，要求读者指出朴素方案会在哪里失败；最后要求读者选择本章机制中的一项或两项做改进。答案不能只给最终代码，还要给不变量、实验矩阵和反例。如果某个答案没有说明失败后如何恢复、如何观察、如何判断优化是否值得，就不能算完整答案。

围绕案例每任务一线程、阻塞混入计算池、移动设备扩展曲线，报告还应补一个迁移问题：同样方法迁移到更大输入、更多线程、不同硬件或本机分布式模拟时，哪些结论保持，哪些结论要重新测。保持的通常是语义边界和不变量；需要重测的通常是吞吐、延迟、批大小、缓存效果、锁竞争和 I/O 行为。这能训练读者区分“原理可以迁移”和“数字不能照搬”。

最后再强调一次本书的写法要求：先从问题出发，再推导机制，再落到实验。不要把实验放成装饰，也不要把 Linux 源码阅读放成资料链接。实验必须检验一个假设，源码阅读必须解释一个用户态现象，项目代码必须保护一个明确边界。读者能按这个顺序复述本章，才说明本章内容真正连贯起来。




\topic{事故复盘式走查}

为了把本章内容真正串起来，可以把把日志作业拆成多个可执行任务写成一次小型事故复盘。事故现象不是一句“系统慢了”或“结果错了”，而要具体到可观察事实：哪一批输入、哪个阶段、哪一个指标、哪一种失败注入触发了问题。如果现象描述不具体，后面的任务和线程、过度订阅和阻塞就只能变成猜测。

复盘第一步是还原朴素设计。写清当时为什么选择它，它解决了什么简单问题，又默认了哪些条件。很多坏设计一开始并不坏，只是从小输入、小并发、无失败的条件迁移到把日志作业拆成多个可执行任务时，没有同步升级边界。这也是本册反复强调从朴素方案出发的原因：只有理解朴素方案为什么成立，才能准确说明它为什么在新条件下失败。

复盘第二步是列出候选假设，但每个假设都必须有可证伪实验。如果怀疑任务和线程相关路径，就构造只改变这一因素的对照；如果怀疑过度订阅，就记录它对应的状态或指标；如果怀疑阻塞，就找出它在代码里的所有入口和退出点。不能把所有怀疑同时改掉。一次修改多个因素会让复盘变成经验叙事，而不是工程证据。

复盘第三步是修复。修复不等于把代码改到通过测试，而是把不变量写回系统。本章的不变量来自“线程数量、任务粒度、阻塞和调度位置没有边界，导致过度并行、上下文切换和不可复现实验”这个失败主题。因此修复说明要写清：新增字段保护什么，新增队列容量限制什么，新增版本号拒绝什么，新增 reference 比较什么，新增指标暴露什么。如果修复无法对应到不变量，就很可能只是局部补丁。

复盘第四步是验收。以“移动设备扩展曲线”为例，验收不应只跑正常输入，还要跑边界输入、压力输入和错误输入。正常输入证明功能还在，边界输入证明不变量没有被破坏，压力输入证明成本模型仍然成立，错误输入证明失败不会越过提交边界。验收报告要保留原始命令和数据，否则下一次回归无法判断问题是否真正修复。

最后要写“未解决问题”。例如调度 trace相关边界可能还没有在当前设备上完全验证，某些 Linux 计数器可能因为权限不可用，某些 NUMA 结论可能因为硬件不支持只能停留在设计层。把未验证风险写出来不是降低质量，而是防止教材把假设写成事实。高质量书籍要敢于说明证据边界，这比把每个结论都写成绝对经验更可靠。

读者完成本章后，可以用这份复盘模板检查自己的学习结果：能否描述事故，能否还原朴素设计，能否提出可证伪假设，能否把修复绑定到不变量，能否设计验收矩阵，能否写出未验证风险。如果六项都能完成，本章知识就已经从概念记忆转成工程判断。如果只能完成其中一两项，就应回到本章前面的案例重新走一遍，而不是急着进入下一章。



\section{本章落到贯穿项目}

Compute Systems Engine 在本章要增加运行环境记录和线程策略实验。第一，程序启动时记录硬件拓扑和可用线程数，至少输出逻辑 CPU 数和用户指定 worker 数。第二，归约 benchmark 要支持从 1 到多倍逻辑 CPU 的线程扫描。第三，实验报告要记录 context-switches 和 migrations。第四，设计线程池前先保留每次创建线程版本作为对照，避免引入线程池后不知道收益来自哪里。

后续项目可以逐步加入亲和性选项，例如不绑定、按物理核心绑定、按 NUMA 节点绑定。每种绑定都必须能失败并报告原因，不能在不支持的平台上默默假装成功。手机或 Termux 环境尤其要谨慎，因为权限、CPU 拓扑和系统调度策略可能和服务器不同。

\topic{线程状态和等待图}

线程在系统里并不只有“运行”和“结束”。它可能 runnable 但还没拿到 CPU，可能 running，可能 blocked 在 futex，可能 sleeping 等 I/O，可能 stopped，可能被 cgroup throttled。分析并发程序时，要把线程状态画进等待图。一个任务慢，可能不是它计算慢，而是它一直没被调度；一个线程池空闲，可能是上游没有提交，也可能是所有任务阻塞在下游队列。

等待图比线程列表更有用。节点可以是线程、队列、锁、I/O 设备、future、远端服务；边表示等待关系。若线程 A 等锁 L，锁 L 被线程 B 持有，线程 B 等队列 Q，队列 Q 等消费者 A，就形成等待环。很多死锁和线程池耗尽都能用等待图解释。无论是 mutex 死锁、future 死锁，还是分布式请求互相等待，本质都是等待关系没有退出路径。

\topic{Runnable 不等于 Running}

一个线程处于 runnable 状态，说明它可以运行，但不代表它正在运行。如果 runnable 线程数长期大于可用 CPU 数，调度器必须轮流运行它们，context switch 增加，每个线程获得的 CPU 时间下降。CPU-bound 程序开远多于核心数的线程，通常不会更快，反而增加调度开销和 cache 污染。

I/O-bound 程序可以有更多线程，因为很多线程在等待 I/O，不同时占用 CPU。但这也有限度：线程栈、调度结构、唤醒风暴和内存占用都会增加。异步 I/O 和事件循环的价值，就是减少“每个等待一个线程”的成本。选择线程数要看任务是 CPU-bound、memory-bound、I/O-bound 还是混合。

\topic{大小核和移动设备}

手机和某些笔记本使用大小核架构，不同核心性能和能效不同。线程绑定到大核和小核，性能可能差很多；系统还会根据温度、电源和前台状态迁移线程。Termux 环境下，用户不一定能完全控制调度。报告中如果在手机上测性能，应记录设备型号、电源状态、温度大致情况、可见 CPU 和是否允许 affinity。

大小核让“线程数等于核心数”更复杂。八个核心不代表八个等价执行资源。CPU 密集任务可能更适合使用大核数量作为主并发度；后台任务可以放小核；长时间满载可能降频，短时间 benchmark 不代表持续性能。第二册的实验要尊重设备现实，而不是假设所有机器都是服务器。

\topic{Oversubscription 的后果}

Oversubscription 指可运行线程数超过可用硬件执行资源。适度 oversubscription 可以隐藏 I/O 等待；过度 oversubscription 会带来上下文切换、cache 失效、锁持有者被抢占、尾延迟增加和内存压力。线程池默认大小不应盲目等于逻辑 CPU 数乘很多倍，除非任务大部分时间阻塞。

更糟的是嵌套并行。上层已经用 8 个线程处理请求，每个请求内部又调用一个默认开 8 线程的库，瞬间产生 64 个 runnable 线程。矩阵库、压缩库、并行 STL 和业务线程池都可能叠加。工程系统应限制嵌套并行，或者让库共享同一个运行时。性能报告要记录是否调用了内部多线程库。

\topic{调度公平和吞吐}

公平调度让每个可运行任务获得合理 CPU 时间，但高吞吐系统有时希望减少迁移和切换。公平、局部性和优先级之间有张力。一个后台压缩线程如果和前台请求线程公平竞争 CPU，可能增加用户请求延迟；完全饿死后台线程，又可能导致 checkpoint 或日志堆积。调度策略要反映业务优先级。

用户态线程池也有类似问题。FIFO 队列公平但可能让长任务阻塞短任务；优先级队列改善前台延迟但可能饿死后台；工作窃取提高负载均衡但可能牺牲数据本地性。操作系统调度和用户态调度叠加，最终行为取决于两者共同作用。

\topic{Run queue 和唤醒延迟}

当一个阻塞线程被唤醒，它不是立即运行，而是进入某个 CPU 的 run queue，等待调度器选择。若系统负载高，唤醒延迟可能明显。条件变量 \code{notify_one} 只是让等待线程有机会运行，不保证它马上持有锁。被唤醒线程还要和其他线程竞争 mutex，可能再次睡眠。

这解释了为什么高竞争条件变量和锁会产生尾延迟。一个生产者通知消费者后，消费者可能晚些才运行；消费者醒来后发现队列又被其他消费者取空；再睡回去。正确性由谓词循环保证，性能则需要减少惊群、控制队列容量、降低竞争和让任务粒度合适。

\topic{优先级反转的系统版}

优先级反转不只存在实时系统。一个低优先级后台任务持有全局锁，高优先级请求线程等待这把锁，而中优先级 CPU 任务不断运行，使后台任务迟迟无法释放锁，这就是普通服务中的优先级反转形态。解决方法包括缩短临界区、避免后台任务持有前台锁、资源隔离、优先级继承或使用无锁/分片结构。

在分布式系统中也有类似现象：低优先级 checkpoint 占满 I/O，导致高优先级提交延迟；后台 compaction 占满 CPU，导致 RPC 超时。线程调度的公平性问题会在资源调度层面反复出现。第二册把这些主题放在同一册，是为了让读者看到共性。

\topic{线程栈和内存占用}

每个线程都有栈。线程数很多时，栈虚拟地址和实际提交内存都可能成为问题。递归、大对象局部变量和深调用链会增加栈压力。C++ 并发代码中，任务对象如果捕获大 vector 按值复制到线程栈或闭包里，也会造成内存和复制成本。线程池减少线程数量，也能减少栈开销。

用户态 fiber 或 coroutine 可以用更轻量的栈或无栈状态机，但它们引入自己的调度和生命周期语义。第二册不展开 coroutine 细节，但要让读者知道：用更多并发单位隐藏等待时，线程不是唯一选择。线程适合 CPU 执行资源，异步任务适合等待很多 I/O。

\topic{信号和异步打断}

Linux 线程还可能收到信号。信号处理函数会在异步位置执行，能做的事情非常有限。高性能和并发代码通常避免在信号处理函数里做复杂逻辑，而是写入 pipe、eventfd 或设置原子标志，让主循环在安全位置处理。信号也可能打断系统调用，导致返回 EINTR。I/O 代码必须考虑这种情况。

这部分不是本册重点，但它提醒读者：真实系统中的控制流不只来自函数调用和线程调度，还可能来自信号、取消、超时和内核回调。状态机要能处理被打断和重试。

\topic{ThreadSanitizer 和调度测试}

ThreadSanitizer 能发现许多数据竞争，但它会改变性能，不能用于性能测量。调度压力测试可以通过增加线程数、插入 yield、随机 sleep、缩小队列容量和重复运行来暴露时序 bug。不要在正式性能报告中开启 TSan，也不要因为 TSan 通过就认为没有并发逻辑错误。

线程调度 bug 常常低概率出现。测试应固定随机 seed，记录失败场景，并尽量构造小模型复现。例如有界队列容量设为 1，比容量很大更容易暴露等待和唤醒错误；线程池 worker 数设为 1 或 2，比 16 更容易暴露嵌套等待死锁。小规模极端测试是并发正确性的朋友。

\topic{线程池大小的配置}

线程池大小应该是配置，不应写死。默认值可以来自硬件并发度，但要允许用户覆盖。CPU-bound 池、I/O 池和后台池应有不同默认策略。若运行在容器中，还要尊重 cpuset 和 CPU 配额。读取 \code{std::thread::hardware_concurrency} 失败或返回不准时，应有保守 fallback。

配置还要进入实验记录。一个 benchmark 只写“默认线程池”，读者不知道默认是多少，也不知道是否与机器有关。项目输出应显示 \code{worker_count}、\code{hardware_concurrency}、绑定策略和是否使用 SMT。这样结果才可解释。

\topic{线程设计的反例}

常见反例包括：每个请求创建一个线程导致高并发时内存爆炸；CPU-bound 任务开数倍于核心的线程导致调度开销增加；线程池 worker 内阻塞等待同池子任务导致死锁；把 I/O 阻塞任务放入计算池导致计算饥饿；在线程局部缓存里保存请求状态导致任务迁移后读不到；全局队列在多 socket 下成为共享热点；忽略关闭语义导致析构时悬挂。

这些反例不是孤立 bug，而是线程、任务、资源和状态边界没有分清。学习线程调度的目标，是在设计阶段就能问出这些问题。

\topic{案例：日志解析线程数}

日志解析看起来 CPU 密集，但它可能同时受 I/O、分支和内存带宽影响。若输入来自 SSD，单线程解析可能跟不上读取；增加线程能提高解析吞吐。若输入来自慢存储，增加解析线程只会让线程等待 I/O。若解析产生大量错误日志，输出 I/O 又可能成为瓶颈。线程数选择必须和全链路一起看。

实验可以设置三种输入源：内存中已加载字符串、page cache 热文件、冷文件或模拟慢读取。对每种输入扫描 worker 数。内存输入主要测解析和内存带宽；热文件混合 page cache；慢读取则测 I/O。若三组结果差异大，说明线程数不能脱离输入源讨论。这个案例能训练读者避免“一律开满核心”的经验主义。

\topic{案例：后台任务干扰}

假设 Compute Systems Engine 在 reduce 阶段同时做 checkpoint。Checkpoint 线程写大文件并 fsync，可能占用 I/O；如果它还做压缩，可能占用 CPU；如果它持有 driver 元数据锁，可能阻塞 commit。前台 reduce 变慢时，不能只看 reduce 代码，还要看后台任务是否干扰。

解决方法包括资源隔离、限速、独立线程池、优先级和调度窗口。Checkpoint 可以在 stage 边界执行，也可以后台增量写；压缩可以限制线程数；元数据锁应只保护状态切换，不覆盖大文件写入。线程调度和 I/O 背压在这里交汇。

\topic{线程池配置实验}

项目应提供线程池配置实验。参数包括 compute worker 数、I/O worker 数、是否允许阻塞任务进入 compute 池、队列容量和任务粒度。实验输入包括 CPU-bound map、I/O-bound read、混合 pipeline。指标包括吞吐、队列长度、context switches、worker idle time、任务延迟和超时。

通过实验，读者可以看到：CPU-bound worker 过多会增加调度成本；I/O-bound 任务放 compute 池会让计算饥饿；队列容量太小会频繁背压，太大会增加延迟；任务粒度太细会增加调度开销。线程池配置不是一次性常量，而是 workload 参数。

\topic{用户态调度和内核调度的边界}

用户态线程池决定哪个任务交给哪个 worker，内核调度器决定哪个线程在哪个 CPU 上运行。用户态运行时看不到所有系统负载，内核调度器不知道用户任务的业务优先级。两者之间存在信息差。绑定 CPU、设置线程优先级、分离线程池、使用 cgroup，都在试图协调这两个调度层。

项目文档应说明哪些由用户态控制，哪些由内核控制。比如 worker id 到 CPU id 的映射是亲和性策略；task 到 worker 的映射是运行时策略；worker 是否被抢占是内核调度；页面在哪个 NUMA 节点是内存策略。把边界说清，排查问题时才知道该看哪一层。

\topic{线程池和请求上下文}

请求上下文不应该隐式绑定到线程。在线程池中，一个请求的 parse、compute、write 可能由不同 worker 执行；即使当前实现碰巧同一线程，后续 work stealing 或异步 I/O 也可能改变。上下文应作为显式对象传递，或者由 task id 在状态表中查找。线程局部只适合 worker 缓冲、统计和临时 scratch。

这个原则能避免很多隐蔽 bug。把 request id 放在线程局部，日志可能串到别的请求；把取消标志放线程局部，后续任务看不到；把输出缓冲放线程局部并在异步写完成前复用，会产生数据损坏。线程局部是性能工具，不是业务上下文容器。

\topic{调度章节验收清单}

本章项目验收包括：能记录硬件并发度和用户 worker 数；能扫描线程数并输出扩展曲线；能记录 context switches 和 migrations；能区分 CPU-bound 和 I/O-bound 任务；线程池配置可调；任务粒度可调；请求上下文显式传递；阻塞任务不会默认进入计算池；报告说明当前环境的 CPU 拓扑限制。满足这些条件，后续线程池和任务图才有可靠基础。

\topic{调度章节总结}

线程是执行资源，任务是逻辑工作，二者不能混淆。调度优化要同时考虑任务粒度、阻塞、亲和性、上下文切换、线程局部缓存和内核调度。线程池能减少线程创建成本，但引入队列、等待、关闭和嵌套提交语义。正确的线程设计不是“开几个线程”，而是让执行资源、数据位置、等待关系和业务上下文都有清楚边界。

\topic{调度决策表}

配置线程和任务时，先问任务是否 CPU-bound、是否阻塞、是否有依赖、是否需要优先级、是否访问大块本地数据、是否会嵌套提交。CPU-bound 任务按核心控制并发；阻塞任务隔离到 I/O 池；有依赖任务进入任务图；有优先级任务需要多队列或配额；访问大块数据要考虑亲和性；嵌套提交要禁止阻塞等待或支持 work helping。这个表能防止线程池被当成万能容器。

\topic{案例：任务粒度自适应的观察}

自适应任务粒度不一定要一开始自动实现，但可以先观察。记录每个任务耗时、队列等待、worker 空闲时间和任务数。如果任务耗时远小于队列操作成本，说明粒度太细；如果 worker 空闲多且少数任务很长，说明粒度太粗或负载不均。运行时可以先把这些指标输出给报告，让人手动调整 cutoff。自动调节只是把这个反馈回路程序化。

这个案例提醒读者：很多“智能调度”可以从简单观测开始。没有观测，自动策略只是猜测；有了观测，即使手动调参，也比固定常数可靠。

\topic{案例：线程池里的阻塞隔离}

假设 compute pool 有 8 个 worker，其中 8 个任务都在等待磁盘写完成，那么新的计算任务无法执行，CPU 可能空闲但队列积压。解决方法不是简单增加 compute worker，而是把阻塞写入交给 I/O pool 或异步写出。Compute worker 只生成输出 block，然后返回执行资源。I/O pool 完成后通过 completion 通知 driver。

这个案例说明线程池大小不是唯一问题，任务类型也必须分类。CPU-bound、I/O-bound、control-plane、background maintenance 应有不同资源策略。否则一个慢资源会拖住所有执行资源。

\topic{运行队列视角下的线程池}

用户态线程池的队列和内核 run queue 是两层不同队列。任务先排在用户态 ready queue 里，worker 线程取到任务后进入执行；worker 自己又作为内核任务排在某个 CPU 的 run queue 里。用户态队列空，不代表系统没有可运行线程；内核 run queue 繁忙，也不代表线程池有任务。分析吞吐和延迟时，必须把这两层分开。

例如一个服务的用户态任务队列持续增长，但 CPU 利用率很低。可能原因不是 worker 数不足，而是所有 worker 阻塞在 I/O 或锁上，内核看来它们不可运行。再例如 CPU 利用率很高，任务队列也增长。可能是 worker 正在运行但每个任务太慢，也可能是系统里其他进程占用 CPU，让 worker runnable 但很少 running。只看一个队列会误判。正确报告应同时记录用户态队列长度、worker 状态、CPU 利用率、context switch、run queue 压力和阻塞原因。

Linux 的调度器不会理解用户态任务优先级。若线程池只有一个 FIFO 队列，短任务和长任务混在一起，内核只看到 worker 都在运行，不知道某个 worker 正在执行低优先级长任务。若业务需要前台请求优先，就要在用户态运行时设计优先级队列、时间片、任务拆分或隔离线程池。操作系统提供 CPU 时间，业务调度要由程序自己表达。

这个视角也解释了 work stealing 的边界。工作窃取能缓解用户态任务分配不均，让空闲 worker 从其他队列拿任务；但如果某些 worker 在内核层面长期得不到 CPU，或者被 cgroup 限额节流，窃取策略无法凭空创造执行资源。运行时优化要先确认瓶颈在用户态队列、内核调度、I/O 还是硬件资源。

\topic{尾延迟和调度抖动}

吞吐系统关注总任务数，低延迟系统还关注尾部任务。一次迁核、一次锁持有者被抢占、一次缺页、一次 cgroup throttling，都可能只影响少数请求，却显著拉高 p99 或 p999。平均耗时变好，不代表用户体验变好。线程调度章节必须把尾延迟作为一等指标，而不是只看总时间。

尾延迟分析要按路径拆开。请求进入队列前等待了多久，排队多久，被哪个 worker 取走，执行中是否阻塞，是否迁移，是否等待锁，是否等待下游 I/O，是否被后台任务干扰。每一段都应有时间戳或事件记录。只有总耗时一个数，就无法知道尾部来自排队、执行、同步还是 I/O。Compute Systems Engine 的任务记录可以逐步加入 submit time、start time、finish time、worker id 和重试次数，这些字段比复杂优化更早产生价值。

调度抖动还有环境因素。手机设备会降频，笔记本会进入省电模式，云服务器可能有邻居噪声，容器可能被 CPU quota 周期性限制。实验报告若不记录环境，尾延迟结论很难复现。对于教学项目，可以先要求读者写出“本次实验不保证独占机器”这样的限制，再逐步学习如何隔离变量。

\topic{亲和性不是固定答案}

亲和性常被误解成“绑定一定更快”。实际上，绑定只是减少调度自由度，让执行位置更可控。若任务的数据也在相同 NUMA 节点，绑定可能提高局部性；若数据在远端节点，绑定会稳定地产生远端访问；若系统有其他干扰，适度不绑定反而让调度器能避开忙 CPU。亲和性是一种实验变量，不是默认优化。

一个合理实验应至少比较四组。第一组完全不绑定，作为系统默认行为。第二组只绑定线程，看缓存热度和迁移是否改善。第三组只控制内存放置，观察线程迁移后远端访问。第四组同时绑定线程和内存，测试最佳局部性。若平台不支持 NUMA，也可以比较不绑定、绑定到物理核心、绑定到 SMT sibling、绑定到一组核心。每组都要记录绑定是否成功，因为容器和手机环境可能拒绝或修改 affinity。

亲和性还要和任务分片配合。若 worker 0 固定在 CPU 0，却不断处理来自所有分片的数据，绑定价值有限。更好的策略是 worker、数据分片和队列分片对齐：某个 NUMA 节点的 worker 优先处理本节点数据，跨节点窃取作为兜底。这样能把亲和性从“钉线程”提升为“数据和执行位置共同设计”。后续分布式章节也会看到同样思想：计算应尽量靠近数据。

\topic{线程池指标的最小闭环}

线程池想要可调，必须先可观测。最小指标包括提交任务数、完成任务数、队列长度、worker 空闲时间、任务排队时间、任务执行时间、阻塞任务数、拒绝任务数、关闭状态和异常任务数。再进一步，可以记录每个 worker 的任务数、最大连续忙碌时间、窃取次数、空轮询次数和最后一次任务类型。没有这些指标，线程池参数只能凭感觉调整。

指标本身不能破坏性能。高频写指标应使用 worker 本地槽位或 relaxed 原子，周期性汇总；低频控制状态可以用锁保护。不要为了监控每个任务而在热路径写全局锁日志。更好的方式是采样、分片、批量写入和可配置开关。监控代码也是并发系统的一部分，它会改变被监控对象。

指标解释也要谨慎。队列长度高可能是任务太多，也可能是 worker 阻塞；worker 空闲高可能是上游供给不足，也可能是任务粒度太粗导致尾部不均；context switch 高可能是 I/O 服务正常现象，也可能是 CPU-bound 程序过度线程化。指标不是结论，而是提出下一步问题的证据。

\section{本章习题}

\begin{exercisebox}
习题一：给当前 \code{parallel_sum} 写一份线程生命周期分析。指出每次调用创建多少线程，线程做多少实际工作，join 等待在哪里发生，哪些成本会在输入很小时占主导。

习题二：设计一个任务粒度实验。固定总输入大小和 worker 数，改变 chunk size。报告要包含任务数、总时间、每个 worker 任务数和尾部完成时间。解释粒度太小和太大分别会出现什么现象。

习题三：构造线程池死锁案例：worker 任务提交子任务并等待。画出等待图，说明为什么增加队列容量不能解决这个死锁。再提出两种运行时语义上的修复方案。

习题四：记录本机或手机上的 CPU 拓扑。写清逻辑 CPU、物理核心或可见核心、是否能看到 SMT、是否有 NUMA 信息、是否允许设置 affinity。不要把无法获取的信息编造成确定结论。
\end{exercisebox}

\begin{keyidea}
线程设计的核心不是“开几个线程”，而是决定哪些状态线程本地，哪些状态共享，哪些共享状态需要同步，哪些等待应该阻塞或异步化。
\end{keyidea}

\begin{labbox}
实验：把并行归约的 worker 数从 1 增加到逻辑 CPU 数的两倍。记录时间、context-switches 和 migrations。解释从哪个点开始扩展性变差。
\end{labbox}
